\documentclass[11pt,a4paper]{article}
\usepackage[T1]{fontenc}
\usepackage[utf8]{inputenc}
\usepackage{lmodern,microtype}
\usepackage[margin=.85in]{geometry}
\usepackage{amsmath,amssymb,amsthm,mathtools,aliascnt}
\usepackage{booktabs,tabularx,array,enumitem,needspace,float}
\usepackage[section]{placeins}
\usepackage{xcolor,tikz,mdframed,graphicx}
\usetikzlibrary{arrows.meta,positioning,calc,decorations.pathreplacing}
\usepackage[hidelinks]{hyperref}
\usepackage[nameinlink,noabbrev]{cleveref}
\usepackage{fancyhdr}
\hypersetup{pdftitle={Sync HotStuff RR, Part II v3.3: Sequencer Rotation},pdfsubject={Espresso-confirmed rotation with automatic and multisignature Force}}
\setlength{\headheight}{14pt}
\pagestyle{fancy}\fancyhf{}
\lhead{\small Sync HotStuff RR, Part II}\rhead{\small Version 3.3 / Reviewed interfaces}
\cfoot{\thepage}
\setlength{\parindent}{0pt}\setlength{\parskip}{4pt}
\setlength{\emergencystretch}{2em}
\renewcommand{\topfraction}{.92}
\renewcommand{\textfraction}{.06}
\renewcommand{\floatpagefraction}{.85}
\makeatletter
\setlength{\@fptop}{0pt}\setlength{\@fpbot}{0pt plus 1fil}
\makeatother
\setlist{nosep,leftmargin=*}
\newcolumntype{Y}{>{\raggedright\arraybackslash}X}
\definecolor{rrRed}{HTML}{A51C30}
\definecolor{rrBlue}{HTML}{2563EB}
\definecolor{rrTeal}{HTML}{0F766E}
\newcommand{\RR}[1]{{\color{rrRed}#1}}
\newcommand{\QC}{\mathcal C}
\newcommand{\FC}{\mathcal F}
\newcommand{\Hash}{\mathsf H}
\newcommand{\rank}{\operatorname{rank}}
\newcommand{\TB}{\mathit{TB}}
\newcommand{\tagg}[1]{\ensuremath{\mathsf{#1}}}
\newcommand{\defterm}[1]{\textcolor{rrBlue}{\textit{#1}}}
\newtheorem{definition}{Definition}[section]
\newaliascnt{assumption}{definition}
\newtheorem{assumption}[assumption]{Assumption}
\aliascntresetthe{assumption}
\crefname{assumption}{assumption}{assumptions}
\Crefname{assumption}{Assumption}{Assumptions}
\newaliascnt{proposition}{definition}
\newtheorem{proposition}[proposition]{Proposition}
\aliascntresetthe{proposition}
\crefname{proposition}{proposition}{propositions}
\Crefname{proposition}{Proposition}{Propositions}
\newcounter{protocol}
% Protocol boxes float to the top of a page, so long boxes leave no gaps.
\newenvironment{protocol}[3][htbp]{\begin{figure}[#1]\refstepcounter{protocol}\label{#2}
\begin{mdframed}[linewidth=.6pt,skipabove=0pt,skipbelow=0pt,
innertopmargin=7pt,innerbottommargin=7pt]\small
\textbf{Protocol \theprotocol. #3}\par\smallskip}{\end{mdframed}\end{figure}}
\crefname{protocol}{Protocol}{Protocols}
\Crefname{protocol}{Protocol}{Protocols}
\title{Sync HotStuff RR, Part II: Sequencer Rotation\\\large Version 3.3: Reviewed Interfaces and Selection Repairs}
\author{}\date{18 October 2026}
\begin{document}
\maketitle
\begin{abstract}
Sync HotStuff RR lets several regions take turns producing one chain.
Each region runs the Part~I engine with a fixed voting snapshot and a
fixed quota of DATA blocks. A terminal regional receipt opens the next
normal session without waiting for global confirmation. A Publish Layer
guarantees submission of complete regional results to Espresso. Espresso
provides bounded finalization and visibility to online readers. An observer confirms a result
globally when it accepts that result as a continuous extension in
Espresso's finalized history. External control uses Espresso timestamps.
A block-missing Force follows from insufficient accepted DATA in a fixed
window. A multisignature Force is announced on L1 before its scheduled
Espresso time. Both preserve the entire global prefix and close the old
production branch to further global output. Each node manages concurrent session programs with explicit preparation,
entry, completion and retirement. Normal operation retains the regional
DATA/DRAIN engine. Authenticated retirement adds a signing and entry gate
while existing timers and evidence processing drain. We specify the
interfaces, event order and conditional safety arguments.
The progress conditions are stated separately. The publication service is
abstract in the protocols; a leader-first implementation with backup
senders and fee sponsorship is described at the end.
\end{abstract}

\section{Overview}
\label{sec:overview}
Part~I~\cite{part1} gives one region a synchronous BFT engine.
This document connects regional engines into one chain. The full view is
\[
 (\beta,r,g,t).
\]
Here $\beta$ is the \defterm{branch version}, $r$ is the regional turn
number, $g$ is the recovery number within that turn, and $t$ is the
leader term within that recovery. A \defterm{session} is the pair
$(\beta,r)$: one region's production turn. Its voting snapshot and DATA
quota stay fixed through all of its recoveries.

A normal regional handoff increases $r$ and keeps $\beta$.
An effective Force increases $\beta$ by one. This number is not an
Espresso validator epoch or a time interval. It distinguishes the old
production branch from the new one. No separate session-ID notation is
needed. Each node and certifier uses these same numbers when replaying
the authenticated history.

\paragraph{What version 3.3 changes.}
Version 3.2 introduced the four-module runtime and six-phase concurrent
program lifecycle. Three fresh reviewers then checked a fixed snapshot.
This revision separates historical handoff validation from permission
to enter now; specifies complete Ready state transfer; and requires
regional multicast of complete receipts to close the Completion case
in the retirement proof. It also advances terminal selection references
and permits same-region renewal when no other region is eligible.
The DATA/DRAIN structure, Force cutoff, old-version check and MSIG/AUTO
priority remain. The entire model section is still identical to v3.1.

These changes do not by themselves prove the full liveness goal.
The raw publication bounds still do not ensure acceptance before each
fixed cutoff. \Cref{sec:liveness-counterexamples} records the stronger
independent liveness traces and their exact premise interpretations.
The Ready timing claim is explicitly conditional on the modeled complete
message handler; arbitrary native installation is not assigned a new
bound. This version remains a candidate for a new independent review.

\paragraph{Protocol outline.}
\begin{itemize}
\item \textbf{Normal rotation remains fast.} A valid receipt for the
terminal DATA block $T_{\beta,r}$ opens the normal successor. The successor
needs the boundary, state and authority, but does not wait for $T_{\beta,r}$
to become globally final.
\item \textbf{Espresso supplies global order.} The Publish Layer submits
complete ranges of regional results. Existing participants implement this
service; it does not run another consensus protocol. Every verifier uses the same
Espresso history and control rules to accept a continuous DATA prefix.
There is no per-block committee acknowledgement or global waiting timer.
\item \textbf{There are two Forces.} An automatic Force follows from
a failed DATA-delivery window. A multisignature Force follows from an
authorized L1 record with a future Espresso activation time.
\item \textbf{A Force preserves all global output.} It closes the old
branch version $\beta$ and starts a new one from the current
global prefix. Normal successors on the old branch version cannot publish its
unconfirmed suffix after closure.
\item \textbf{Control is interpreted by verifiers.} Espresso stores
ordered bytes. Nodes and certifiers verify L1 and Espresso evidence and
run the same external interpreter. Espresso need not execute a contract
for this chain.
\end{itemize}

The regional engine retains DATA, recovery-only DRAIN, FINAL receipts,
and Completion. A DRAIN has no application payload and never enters the
global DATA prefix. Force closure does not cancel a regional timer,
change an existing lock or fabricate a quit proof. Its separate retirement
request stops new work and drains existing work under \cref{sec:retirement}.

\begin{table}[htbp]\centering\small
\caption{Participants and responsibilities.}\label{tab:roles}
\begin{tabularx}{\linewidth}{@{}>{\raggedright\arraybackslash}p{.19\linewidth}Y@{}}\toprule
Role & Responsibility\\\midrule
L1 & Publishes immutable authorities and the ordered control log.\\
Regional replica & Runs the DATA/DRAIN engine of its sessions and
serves complete commitment evidence.\\
Multisignature service & Announces future status changes and Forces
on L1. It cannot choose a retained prefix or declare data global.\\
Publish Layer & Delivers complete regional results to Espresso. The
leader, replicas and external verifiers may act as senders.\\
Espresso & Supplies a finalized ordered log, authenticated timestamps,
complete namespace contents, and bounded finalization and visibility.\\
Observer / certifier & Verifies both sources, replays control and
exports, and derives global output and forced boundaries.\\
\bottomrule
\end{tabularx}
\end{table}

\paragraph{Production, confirmation, and publication.}
The \defterm{production branch} $R_u$ of observer $u$ may contain
several normal sessions beyond its \defterm{global prefix} $G_u$.
The \defterm{export session} is the session responsible for the next
DATA height after $G_u$. It need not be the session currently producing
new blocks. The accepted DA publication prefix $C_{\rm pub}$ may lag
behind $G_u$ as well. Neither lag permits a Force to remove $G_u$.
The Publish Layer delivers results to Espresso; $C_{\rm pub}$ instead
tracks the later rollup DA publication stage.

A block's regional receipt remains evidence of its regional commitment
even if its unconfirmed branch is later closed. Global finality is
stronger: accepted DATA is never removed by either Force.
The claim concerns input content, order and provenance. Execution
correctness and native state continuity remain obligations of the
application and publication interfaces.

\begin{table}[htbp]\centering\small
\caption{Changes of production view.}\label{tab:views}
\begin{tabularx}{\linewidth}{@{}>{\raggedright\arraybackslash}p{.19\linewidth}>{\raggedright\arraybackslash}p{.25\linewidth}Y@{}}\toprule
Change & New view & Trigger and timer effect\\\midrule
Normal leader rotation & $(\beta,r,g,t+1)$ & A terminal QC below $H_{\beta,r}$;
existing commit timers continue.\\
Regional recovery & $(\beta,r,g+1,0)$ & A quit proof, the $\Delta_i$ wait
and finalized authority; the quit cancels generation-$g$ timers.\\
Normal regional handoff & $(\beta,r+1,0,0)$ & Ready on a REGULAR witness
with the terminal receipt; the old engine stops through Completion.\\
Forced handoff & $(\beta+1,r_c+1,0,0)$ & Ready on a derived Force boundary;
$r_c$ is the boundary's predecessor index. Old programs close separately.\\
\bottomrule\end{tabularx}
\end{table}

A Force can reuse an index from a discarded normal successor.
The new pair $(\beta+1,r_c+1)$ keeps votes, locks and evidence separate.
For example, $(7,20,0,0)\to(7,21,0,0)$ is a normal regional
handoff. A Force that closes export session $(7,20)$ instead opens
$(8,21,0,0)$, even if an unconfirmed normal successor already used
$(7,21)$. The retained DATA anchor determines where the new session starts.
A finalized but not yet effective L1 control does not stop production
or global acceptance.

\section{Model and assumptions}
\label{sec:model}
\subsection{L1 control state}
\label{sec:l1}
Nodes authenticate finalized L1 headers, authorities and a complete
ordered control prefix. L1 timestamps label L1 blocks; they are not
regional commit times or Espresso activation times.

\begin{table}[htbp]\centering\small
\caption{L1 authority and control records.}\label{tab:l1}
\begin{tabularx}{\linewidth}{@{}p{.27\linewidth}Y@{}}\toprule
Field & Meaning\\\midrule
$\tagg{AUTH}_{r,0}$ & The fixed regional snapshots, leader schedule,
DATA quota $\TB_r$, and export-window parameters for index $r$.\\
$\tagg{AUTH}_{r,g}$, $g>0$ & The leader schedule for recovery
generation $g$; it does not change the original snapshot or quota.\\
Control sequence $c$ & A unique increasing index in the L1 control
log. It binds the chain, record kind and payload.\\
Activation time $\theta_a$ & An Espresso timestamp threshold,
announced before it can be crossed.\\
STATUS & A future live/off update for the named regions. It changes
selection eligibility, not existing voting weights.\\
MSIG FORCE & A target branch version, activation time and any authorized
status changes to apply with the Force. No retained height is supplied.\\
CLEAR & An explicit future removal of the named automatic cooldown;
ordinary ON does not perform this operation.\\\bottomrule
\end{tabularx}
\end{table}

\paragraph{Authority records.}
For each $(r,g)$, L1 supplies at most one immutable authority.
$\tagg{AUTH}_{r,0}$ fixes $\mathit{cfg}_r$, including the members,
weights, $F_i$, $\Delta_i$ and $K$ of each region, and fixes $\TB_r$.
It also fixes positive window lengths $W_{\rm start},W_{\rm run}$,
positive count targets $m_{\rm start},q_{\min}$, and cooldown
$C_{\rm auto}>0$. Counts are capped by the remaining DATA quota.
For $g>0$, only the leader schedule changes.
Nodes wait for required authority evidence. Later controls and slashing
never reweight an old certificate or reset a session quota.

Authority is indexed by $r$, not by a particular speculative branch.
A new branch version that reuses $r$ uses the same authority at that index.
Its pair $(\beta,r)$ identifies a different engine. Rebinding an already
assigned authority is not permitted.

\paragraph{Control ingestion and activation.}
The complete L1 control prefix has a sequence number and a binding
history commitment. A verifier retains pending records as well as
already applied records. Ingestion does not imply activation.
Activation times need not follow sequence order: records becoming due
together are ordered by $(\theta_a,c)$. Malformed, unauthorized or
duplicate records do not grant a second action.

STATUS and CLEAR are independent controls. Configuration attached to
a Force applies only if that Force's target branch version is still current.
A stale Force is consumed without changing the new branch version. An independent
STATUS beside it still takes effect on its own terms.
Records are immutable after finalization. A revision is a new,
prospectively scheduled record, not an edit to past history.

\paragraph{Espresso positions and effective status.}
Let $J_q$ be Espresso block $q$, with authenticated timestamp
$\theta_q$. Positions also include namespace-transaction and record
indices. Heights and those indices give a total order even when
timestamps are equal. For a threshold $d$, define
\[
 q(d)=\min\{q:\theta_q\ge d\}.
\]
Its event is applied at the opening of $J_{q(d)}$, before that block's
EXPORTs. The prelude applies due independent status changes, then a
matching multisignature Force, then any automatic Force.

Write $\chi_E(q)$ for the eligibility snapshot after this prelude and
before the EXPORTs. It combines the L1 live/off map with automatic
cooldowns through that position. A region is eligible only if its
governance flag is on, its cooldown has expired, and its designated
authority has positive total weight. The historical snapshot is
authenticated by replay, not by an endpoint's statement.

OFF changes subsequent selections made from the affected snapshots.
It does not revoke an existing historical REGULAR witness. ON does not
clear cooldown or reopen an old branch version. A committee that needs to stop
the current branch uses a Force. No control changes an accepted DATA
prefix or retroactively changes a regional voting snapshot.

\paragraph{Genesis.}
The initial checkpoint fixes the chain identifier, cyclic order
$\mathcal L_R$, genesis DATA anchor, first replay position, Espresso
quorum trust anchor, control prefix and initial branch version $\beta=0$.
A synthetic boundary selects the first eligible region in
$\mathcal L_R$ at the first eligible authenticated position. It waits
if none exists. The initial authority and state are trusted or verified
from their specified sources; Espresso does not authenticate earlier
application history merely by starting this interpreter.
The genesis anchor includes synthetic commitment evidence and fixed,
authenticated L1 and Espresso references. Its pending boundary has
source context $(i_b,\beta_b,r_b)=(\bot,\bot,-1)$. It opens
session $(\beta,0)$ under the version current at selection: normally
$(0,0)$, or a later version if a Force acted while genesis was Waiting.
At selection height $q_s$, set $q^{\rm open}_{\beta,0}=q_s$ and
$\lambda^{\rm open}_{\beta,0}=\max\{\lambda(\text{genesis}),
\lambda(\tagg{AUTH}_{0,0})\}$, where the second term is the authority's
L1 publication height. A supplied later checkpoint must instead include
its current boundary, reference bounds and complete interpreter state.

\subsection{Data delivery and the Publish Layer}
\label{sec:delivery}
The \defterm{Publish Layer} is a delivery service implemented by existing
participants. It does not run another consensus protocol or choose which
DATA is valid. A regional leader, another replica, or an external verifier
can perform a submission. The rotation rules use the service guarantee
below; \cref{sec:publish-implementation} describes one implementation.

A \defterm{complete regional result} is DATA with its direct or covering
FINAL receipt and the ancestry, boundaries and source evidence needed to
validate it. Correct replicas offer complete results to the Publish Layer
as soon as they obtain them. The service may batch results and construct
EXPORT records as specified in \cref{sec:exports}. A proposal QC, a block
hash, or a receipt whose payload is missing is not a complete input.

Regions may send the same results directly to other regions. This lets a
regular successor reach Ready before the terminal result is finalized in
Espresso. If direct delivery is unavailable, it can obtain the result from
Espresso after the record is finalized. The latter path waits for Espresso; the
protocol does not impose that wait when direct delivery succeeds.
There is no separate cross-region transport service in this model.

Espresso carries the finalized records, including DATA and commitment
evidence, to online readers. Readers still check the application and
rotation rules. L1 supplies its own authority and control proofs.
An execution-state snapshot may be transferred directly or reconstructed
from authenticated input history. State reconstruction is part of Ready,
not something that Espresso inclusion performs automatically.

\subsection{Stake and faults}
\label{sec:stake}
Member $v$ of region $i$ has voting weight $\omega_{i,v}>0$. Let
$W_i=\sum_{v\in V_i}\omega_{i,v}$. The stake configuration
fixes a bound $F_i<W_i/2$, and Byzantine members of region $i$ hold
total weight at most $F_i$. Every regional certificate needs
matching signatures from distinct members with total weight greater than
$F_i$. So every certificate has a correct signer, and the correct members
alone can form one. This replaces the $f+1$ of $n=2f+1$ in Part~I.

Membership and weights stay fixed for a whole session, including its
recoveries. Slashing and status updates do not change them. A later
session checks old certificates against the old snapshot. The bound
$F_i$ counts every key of the snapshot that is compromised while its
certificates are still accepted. Compromise of old keys after that is
outside this model.

Correct replicas keep their protocol state and committed blocks.
Signatures cannot be forged, and hashes are binding.



\subsection{Assumptions}
\label{sec:assumptions}
\begin{assumption}[Finalized L1 records]
\label{ass:authority}
At every finalized L1 height, the authority records and ordered control
history through that height are common and authenticated. Finalized
records never change. A verifier obtains a complete control prefix,
not an arbitrary subset of valid events.
\end{assumption}

\begin{assumption}[Authenticated synchronous communication and global clock]
\label{ass:network}
\leavevmode\par
\begin{itemize}
\item Messages use authenticated communication channels.
\item Within region $i$, the actual delay of a correct-to-correct
message is $\delta\le\Delta_i$; $\delta$ may differ across messages.
\item All nodes share a global clock.
\end{itemize}
\end{assumption}
The delay bound includes delivery and processing of complete validation
dependencies. Each replica records its own entry and voting events on
the global clock.

\begin{assumption}[Bounded publication service]
\label{ass:publish}
If a correct participant has a complete regional result available to the
Publish Layer at real time $s$, the service submits an Espresso-admissible
record carrying that result by $s+T_{\rm pub}$. The service preserves its
DATA and commitment evidence and does not depend on the original leader
alone. At least one working submission path and the required fee resources
are available. Here $T_{\rm pub}>0$ is fixed for the stated input-size and load
limits. A result already recorded in Espresso satisfies this delivery
guarantee; retries needed for acceptance into $G$ are a separate task.
\end{assumption}
The clock starts when the complete result is available, not when a leader
first proposes a block or one replica first commits it. Forming the FINAL
receipt remains a regional-protocol task. ``Submits'' means delivery of the
complete record to the Espresso submission interface under its admission
rules, not merely a local send attempt.


\begin{assumption}[Honest multisignature committee]
\label{ass:committee}
The committee is one honest, coordinated control service. It submits
only authorized, immutable, future-dated STATUS, CLEAR and MSIG FORCE
records. It chooses activation times satisfying the advance-announcement
condition below. It keeps no private protected prefix and has no power
to select an older retained anchor.
\end{assumption}

\begin{assumption}[Bounded L1 finalization and visibility]
\label{ass:monitor}
A valid control operation actually submitted at real time $s$ is
finalized, available and verifiable by every correct online participant
by $s+T_{\rm L1}$. Participants include regional replicas, external
observers and certifiers. The bound includes inclusion, finalization,
delivery, retrieval and validation.
\end{assumption}

\Needspace{8\baselineskip}
\begin{assumption}[Espresso history and time]
\label{ass:espresso}
Correct verifiers obtain compatible finalized Espresso prefixes and
can authenticate each required header and complete namespace payload,
or its genuine absence. The history and evidence needed for replay remain
retrievable. In particular, globally accepted DATA remains recoverable
without the original producing region staying online. Timestamps are
non-decreasing. If block $J_q$ first becomes final at real time $t_E(q)$,
then $\theta_q\le t_E(q)+\epsilon_E$, where $\epsilon_E\ge0$ is the
configured maximum clock lead.
\end{assumption}
Long-term retrieval is an explicit service assumption. It may be supplied
by archival readers that preserve authenticated Espresso history; it does
not follow from a finality certificate alone. A returning node fetches
missing history or starts from an authorized checkpoint before confirming
new output. Catching up with arbitrarily long history is not assigned a
constant-time bound.

\begin{assumption}[Bounded Espresso finalization and visibility]
\label{ass:espresso-visibility}
An admissible complete record submitted to Espresso at real time $s_E$
is included in a finalized block by $s_E+T_E$, for a fixed $T_E>0$.
By the same deadline, every
correct participant that stays online and keeps up with the history can
retrieve the complete record and verify its inclusion and finality.
The bound covers inclusion, finalization, dissemination, retrieval and
verification of the Espresso evidence, under the declared admission and
load conditions.

Each newly finalized Espresso block and its complete relevant namespace
also become available and verifiable to those participants within $T_E$
of finalization, including blocks with no records for this chain.
\end{assumption}
The service provides full payloads and evidence, not only header hashes.
It does not promise that an EXPORT passes this chain's acceptance rules:
a stale base, closed branch or missing predecessor can still prevent it
from advancing $G$. Regional and application validation remain separate.

\paragraph{The two delivery bounds.}
For a newly available complete result, the two assumptions give
\[
 s_{\rm submit}\le s+T_{\rm pub},\qquad
 s_{\rm visible}\le s+T_{\rm pub}+T_E.
\]
Here $s_{\rm visible}$ means that the submitted result and its Espresso
evidence are available to the online readers above. It is not a bound on
when the regional result first forms, when a new region restores its
execution state, or when a rejected EXPORT is accepted after retries.
These are real-time bounds. Delivery windows use Espresso timestamps, so
window parameters also need the clock profile and regional progress and
startup allowances. The bound $T_E$ does not replace the clock-lead
assumption or the advance-announcement rule below.


\paragraph{Advance announcement.}
For a control submitted at $s$, the committee chooses
\[
 \theta_a>s+T_{\rm L1}+\epsilon_E+\eta,
 \qquad \eta>0.
\]
Thus the control is available to correct participants before its
first crossing can become final. This is a model condition, not a
conclusion from an average L1 or Espresso block time. The L1 service
bound already includes dissemination; no second announcement delay is
added. Offline or unsynchronized participants catch up before confirming
new output. Missing control history pauses verification.

\paragraph{Verification interfaces.}
Regional receipt validity and the application-validity checks below
are assumed sound under the stated regional model. A deployed
certifier must authenticate both source histories and bind the control
version used. A signature or TEE alone does not establish that an old
control snapshot is complete. The publication interface additionally
requires sound execution and DA correspondence checks.

\subsection{Conditions for progress}
\label{sec:progress}
The following conditions are separate from prefix safety.
\begin{itemize}
\item Correct leaders receive repeated turns that are actually entered,
including recovery initializers and DATA-producing turns across
sessions. Exposed equivocators are eventually excluded permanently
from later leader schedules.
\item Espresso continues to finalize retrievable history with growing
timestamps and includes this chain's valid exports.
\item Complete results become available to the Publish Layer. Its
submission paths and Espresso meet their stated bounds. Retried EXPORTs
for the next required DATA are accepted in time during some sustained
service interval.
\item Required authorities and recoverable native state become
available. A waiting boundary eventually has another eligible region.
The committee leaves working regions enough uninterrupted service time.
\item Correct clients' admissible requests reach the regions that
continue serving. Requests fit in DATA batches and have only finitely
many older requests ahead of them in each local queue.
\end{itemize}
The bounds begin with a complete result and do not bound regional
recovery, state restoration, or repeated rejected EXPORTs. Window success
therefore still needs the stated serving interval and a compatible
Espresso clock profile. Otherwise repeated automatic Forces can prevent
progress without revoking any global DATA. Espresso failure also stops external time and
global progress. Switching the global finality service is outside this
protocol.

\section{Regional engine: definitions}
\label{sec:regional}
\Cref{tab:notation} lists the principal objects. Regional anchors and
view times use the same four indices as the view.
\begin{table}[htbp]\centering\small
\caption{Notation.}\label{tab:notation}
\begin{tabularx}{\linewidth}{@{}p{.32\linewidth}Y@{}}\toprule
Symbol & Meaning\\\midrule
$(\beta,r,g,t)$; $j$ & Branch version, regional turn, recovery, leader term; fresh position.\\
$(\beta,r)$ & Session: one region\textquotesingle s turn in one branch version.\\
$i_{\beta,r}$; $V_i$ & Session region; its regional members.\\
$\omega_{i,v},W_i,F_i$ & Member weight, total weight, and Byzantine bound $F_i<W_i/2$.\\
$\Delta_i$; $K$ & Regional synchronous delay bound; maximum DATA per term.\\
$\tagg{AUTH}_{r,g}$; $\mathit{cfg}_r$ & Immutable authority; fixed session snapshot.\\
$A_{\beta,r,g,t}$ & Anchor before the first fresh position of that view.\\
$h(B)$; $k_{\beta,r}(B)$ & DATA height; session physical position, including DRAIN.\\
$\TB_r,H_{\beta,r},T_{\beta,r}$ & Session DATA quota, endpoint height, and terminal DATA.\\
$m_{\beta,r,g,t}$; $L(\beta,r,g,t)$ & Fresh term positions; designated leader.\\
$\QC_{\beta,r,g,t}(B)$; ANCHOR QC & Proposal certificate; opening re-certification.\\
$\FC_{\beta,r}(B)$ & Regional FINAL receipt, or its stated same-session coverage.\\
$Q_{\beta,r}^0,Q_v,Q^*$ & Inherited commitment evidence; entry lock; selected justification.\\
$e_v(\beta,r,g,t),\tau_v(P),a_v$ & Real-time entry, validated vote and quit events.\\
$\rho,\lambda(B)$ & Regional finalized L1 reference and its L1 height.\\
$\xi(B),\vartheta(B)$ & Signed Espresso selection reference and its timestamp.\\
$J_q,\theta_q$ & Espresso block and authenticated timestamp.\\
$\chi_E(q),\mathcal L_R$ & Historical pre-payload eligibility snapshot; region cycle.\\
$q_s,q_{\beta,r}^{\rm open}$ & Selection Espresso height; session reference lower bound.\\
$\lambda_{\beta,r}^{\rm open}$ & Separate L1 reference lower bound for a new session.\\
$G_u,R_u,C_{\rm pub}$ & Global DATA prefix; provisional branch; accepted publication prefix.\\
$T_{\rm L1},\epsilon_E$ & Full L1 announcement-service bound; maximum Espresso clock lead.\\
$T_{\rm pub},T_E$ & Complete-result submission bound; Espresso finalization and visibility bound.\\
$W_{\rm start},m_{\rm start}$ & Startup-window length and DATA target.\\
$W_{\rm run},q_{\min},C_{\rm auto}$ & Running-window length, target and cooldown.\\
$E$; Ready & Regional entry proof; full boundary/state/authority validation.\\
\bottomrule\end{tabularx}
\end{table}

\subsection{Blocks, positions, and L1 references}
\label{sec:blocks}
The \defterm{anchor} $A_{\beta,r,g,t}$ is the block immediately before the first fresh
position of view $(\beta,r,g,t)$. We always write all four indices.
\begin{itemize}
\item $A_{\beta,r,0,0}$ is the DATA block inherited from the preceding
session; the first session inherits genesis. It keeps the regional commitment evidence from the session that
produced it: a direct or covering receipt as defined in
\cref{sec:coverage}. Genesis has synthetic evidence.
\item For $g>0$, NEWVIEW chooses the \defterm{recovery anchor} $A_{\beta,r,g,0}$
(\cref{alg:recovery}). It is certified but may be uncommitted.
\item For $t>0$, $A_{\beta,r,g,t}$ is the terminal DATA block of term
$(\beta,r,g,t-1)$, certified by that term's terminal QC.
\end{itemize}
The anchor of a view never changes. The generation starts at
$A_{\beta,r,g,0}$; later terms use their own anchors $A_{\beta,r,g,t}$. ANCHOR voting
re-certifies only $A_{\beta,r,g,0}$, once per generation.

Each block has two positions.
\begin{itemize}
\item The \defterm{DATA height} $h(B)$ counts the DATA blocks on the chain
from genesis to $B$.
\item The \defterm{physical position} $k_{\beta,r}(B)$ counts DATA and DRAIN blocks
within session $(\beta,r)$, with $k_{\beta,r}(A_{\beta,r,0,0})=0$. This local count does not change
the signed fields of $A_{\beta,r,0,0}$.
\end{itemize}
A new block of session $(\beta,r)$ has the form
\[
 B=(\beta,r,k,h,\mathit{type},x,\Hash(B_{\rm parent}),\rho,\xi),
 \qquad \mathit{type}\in\{\tagg{DATA},\tagg{DRAIN}\},
\]
with $k=k_{\beta,r}(B_{\rm parent})+1$. A \defterm{DATA block} carries a batch $x$ of
requests and has $h=h(B_{\rm parent})+1$. A \defterm{DRAIN block} carries no batch
and has $h=h(B_{\rm parent})$.

\paragraph{L1 references.}
The \defterm{reference} $\rho$ names a finalized L1 block. It holds that L1
block's height $\lambda(B)$ and hash, and a registry proof for the L1
control records up to that height. These rules hold:
\begin{itemize}
\item along DATA ancestry, $\lambda$ never decreases;
\item the first DATA block of session $(\beta,r)$ has $\lambda\ge
\lambda^{\rm open}_{\beta,r}$, the L1 lower bound fixed by its boundary
(\cref{sec:witness}; genesis uses \cref{sec:l1});
\item a DRAIN block copies the reference of $T_{\beta,r}$.
\end{itemize}
The authenticated timestamp of the L1 block named by $\rho$ can be
read from its header. It is not a commit timestamp for $B$. Recovery
may commit $B$ much later (\cref{sec:late-receipt}). The
regional Vote step applies no receive-time freshness filter to $\rho$.

The L1 reference retains its regional validation role. It does not
choose the external Force time or the new Espresso selection snapshot.
Regional validity uses the fixed session authority. A later status update
does not change an old receipt's voting weights.

\paragraph{Espresso selection references.}
The additional outer field $\xi(B)$ names a finalized Espresso header
$J_q$ by height and hash, with its authenticated pre-payload eligibility
snapshot. Write $\vartheta(B)=\theta_q$. Its height does not decrease
along DATA ancestry, and the first DATA of session $(\beta,r)$ has
$q\ge q_{\beta,r}^{\rm open}$. The terminal DATA additionally requires
$\operatorname{height}(\xi(T_{\beta,r}))>q_{\beta,r}^{\rm open}$.
DRAIN copies the selection reference of $T_{\beta,r}$,
just as it copies its L1 reference. These references and the session
context are bound by the proposal hash and signatures. They do not
change the native mini-block encoding.

A verifier authenticates the referenced history, directly or from its
validated cache. Global acceptance additionally requires the reference
to precede the containing EXPORT's Espresso block. No reference can
justify itself through that same EXPORT. Historical references carry no
receive-time freshness claim. The reference of $T_{\beta,r}$ starts the normal
successor search in \cref{sec:witness}.

\paragraph{Why the selection reference is needed.}
The terminal reference fixes a common starting point for choosing the next
region. It lets nodes make that choice before the terminal DATA itself is
published to Espresso. It is not proof that this DATA is globally final.

For example, let the region order be A, B, C. At Espresso block $J_{500}$,
B and C are eligible, and A's opening reference is below 500. An announced
L1 status update makes B ineligible at
$J_{505}$. A's terminal DATA 110 binds $\xi(110)=(500,\Hash(J_{500}))$.
One node receives its regional receipt when its Espresso view ends at 503;
another receives it when that view ends at 506. Choosing from each node's
latest local status would give B and C respectively. With the signed
reference, both start at 500 and select B, provided the branch has not
been closed by Force. The historical status is derived from L1 controls
at that Espresso position; Espresso does not execute the status rule.

The later OFF does not revise this regular boundary. The reference gives
a common choice, not a guarantee that the status is fresh. Force remains
the way to close an already authorized branch. Bounded visibility does
not remove the need for a common reference: honest nodes can receive a
receipt at different times. This revision changes only the terminal
lower-bound check: its reference must strictly advance the session's
opening position. A selected program fetches a qualifying finalized
header before Ready, so normal DATA production need not wait for it.
The terminal DATA itself still need not be globally confirmed for normal
handoff. DRAIN retains the terminal DATA's reference.

\subsection{Session quota, term positions, and DRAIN}
\label{sec:quota}
\paragraph{1. Fix the session endpoint.}
Session $(\beta,r)$ may add $\TB_r\ge1$ DATA blocks after $A_{\beta,r,0,0}$. Its \defterm{last
DATA height} is
\[
 H_{\beta,r}=h(A_{\beta,r,0,0})+\TB_r.
\]
The DATA block at this height is $T_{\beta,r}$. Recovery never changes $H_{\beta,r}$.
The quota counts the DATA blocks after $A_{\beta,r,0,0}$ in a proposal's
ancestry, whether committed or not. DRAIN, resends, and ANCHOR votes
use no quota.

\paragraph{2. Produce DATA while space remains.}
If $h(A_{\beta,r,g,0})<H_{\beta,r}$, generation $(\beta,r,g)$ is a \defterm{DATA generation}.
Each term begins at its anchor $A_{\beta,r,g,t}$ and has
\[
 m_{\beta,r,g,t}=\min\{K,\;H_{\beta,r}-h(A_{\beta,r,g,t})\}
\]
new positions, numbered $j=0,\ldots,m_{\beta,r,g,t}-1$. At position $j$, the
DATA block has
\[
 h=h(A_{\beta,r,g,t})+j+1,\qquad
 k=k_{\beta,r}(A_{\beta,r,g,t})+j+1.
\]
A new position is called a \defterm{fresh position}; its block is a
\defterm{fresh block}. The anchor itself is not fresh. The last position
$j=m_{\beta,r,g,t}-1$ is the \defterm{terminal position}, and its QC is the
\defterm{terminal QC}.

If the terminal DATA height is below $H_{\beta,r}$, its QC admits term
$(\beta,r,g,t+1)$, with that terminal block as $A_{\beta,r,g,t+1}$. If it is $H_{\beta,r}$,
there is no next term. Only the last DATA term can have fewer than $K$
positions. The initial generation is always a DATA generation, since
$\TB_r\ge1$.

\paragraph{3. Use DRAIN when recovery starts at the endpoint.}
If $g>0$ and $h(A_{\beta,r,g,0})=H_{\beta,r}$, generation $(\beta,r,g)$ is a
\defterm{drain generation}. It has only term $t=0$ and position $j=0$,
with $m_{\beta,r,g,0}=1$. Its one fresh block is DRAIN, whose parent is
$A_{\beta,r,g,0}$ and whose justification is the generation's ANCHOR QC.
It carries no requests, keeps DATA height $H_{\beta,r}$, and has no next term.

DRAIN supplies a commit timer. ANCHOR voting starts no timer, so
re-certifying $A_{\beta,r,g,0}$ alone does not commit it. At height $H_{\beta,r}$, a
DATA child would exceed the quota. The DRAIN child's timer instead
commits its ancestors, including $T_{\beta,r}$.
\begin{itemize}
\item A DATA generation that reaches $H_{\beta,r}$ does not append DRAIN.
It waits for $T_{\beta,r}$'s timer or recovers after a quit.
\item A later drain generation may start from an earlier DRAIN.
Each interrupted recovery may add one DRAIN, so their total number
is not bounded.
\item The choice depends on the recovery anchor's DATA height, not on
whether a replica has already committed $T_{\beta,r}$.
\end{itemize}
DRAIN stays in the regional log. The external chain contains only DATA,
and the next session inherits $T_{\beta,r}$.

\paragraph{Example.}
Let $h(A_{\beta,r,0,0})=100$, $\TB_r=10$, and $K=4$, so $H_{\beta,r}=110$.
\begin{center}\small
\begin{tabular}{@{}lll@{}}\toprule
View & Anchor height & Fresh DATA heights\\\midrule
$(\beta,r,0,0)$ & 100 & 101--104\\
$(\beta,r,0,1)$ & 104 & 105--108\\
$(\beta,r,0,2)$ & 108 & 109--110\\\bottomrule
\end{tabular}
\end{center}
Recovery from DATA 107 starts a new generation whose first term produces
108--110. Recovery from certified DATA 110 starts a drain generation.
In both cases, the endpoint remains 110.

\subsection{Certificates and signed messages}
\label{sec:certificates}
\Cref{tab:certificates} lists the certificates of a session. Each needs
matching signatures from distinct members of $V_{i_{\beta,r}}$ with total weight
$>F_{i_{\beta,r}}$.

\begin{table}[htbp]\centering\small
\caption{Certificates of a session.}
\label{tab:certificates}
\begin{tabularx}{\linewidth}{@{}>{\raggedright\arraybackslash}p{.25\linewidth}>{\raggedright\arraybackslash}p{.24\linewidth}Y@{}}
\toprule
Certificate & Signed message & Purpose\\
\midrule
\defterm{Proposal QC} $\QC_{\beta,r,g,t}(B)$ & DATA-VOTE or DRAIN-VOTE & Justifies a
child. A terminal QC below $H_{\beta,r}$ also admits the next term.\\
\defterm{ANCHOR QC} & ANCHOR-VOTE & Re-certifies the opening $A_{\beta,r,g,0}$ in
generation $g$, before any fresh proposal.\\
\defterm{BLAME proof} & BLAME & Freezes a generation (\cref{sec:quit}).\\
\defterm{Regional receipt} $\FC_{\beta,r}(B)$ & FINAL & Proves that a correct replica
committed DATA block $B$.\\
\bottomrule
\end{tabularx}
\end{table}

A correct replica casts at most one fresh vote per $(\beta,r,g,t,j)$ and one
ANCHOR vote per $(\beta,r,g)$. DATA and DRAIN votes follow the same rule but
are different signed types. No QC and no ANCHOR vote starts a commit
timer.

\paragraph{Signed context.}
Every signature binds the chain identifier, protocol domain, session
$(\beta,r)$, voting configuration, assigned region and message type.
View messages also bind $(g,t)$ and the L1 authority reference
$\tagg{AUTH}_{r,g}$. Their full view is $(\beta,r,g,t)$.
The other signed fields are:
\begin{itemize}
\item \textbf{PROPOSE and fresh votes:} the generation's opening,
position $j$, block type, physical position, DATA height and block hash;
\item \textbf{NEWVIEW and ANCHOR-VOTE:} the chosen opening. NEWVIEW
also binds its justification and entry proof;
\item \textbf{STATUS:} the sender's entry lock and entry proof;
\item \textbf{BLAME:} the view only, with no opening, block hash or
progress count. All complaints for a view sign the same statement;
\item \textbf{FINAL:} the session and DATA fields below, with no $g$ or $t$.
\end{itemize}
Displayed messages include $\beta$ and $r$ but omit the common chain,
protocol, region and configuration fields where they are unambiguous.
A verifier also checks the boundary that assigned the pair $(\beta,r)$;
the pair alone is not evidence of authority.

\paragraph{Rank and the base justification.}
Within session $(\beta,r)$, a QC formed in term $(\beta,r,g,t)$ for block $B$ has \defterm{rank}
$(g,t,k_{\beta,r}(B))$, compared lexicographically. The ANCHOR QC of generation
$g$ has rank $(g,0,k_{\beta,r}(A_{\beta,r,g,0}))$, just below the generation's first
fresh QC. A fixed total order on certified content breaks ties. Signer
sets and encodings do not matter. The \defterm{opening
identity} of generation $g$ is
$(g,k_{\beta,r}(A_{\beta,r,g,0}),h(A_{\beta,r,g,0}),\Hash(A_{\beta,r,g,0}))$. In this document,
``opening'' always means the generation's opening block $A_{\beta,r,g,0}$, as
in Part~I. QCs are ranked only within the same session $(\beta,r)$.

The anchor $A_{\beta,r,0,0}$ also has a \defterm{base justification} $Q_{\beta,r}^0$: the
anchor together with its direct or covering commitment evidence in the
boundary witness (\cref{sec:coverage,sec:witness}). Its justified block
is exactly $A_{\beta,r,0,0}$, even if the evidence uses a receipt for a higher
old-session block. That higher block is evidence only. $Q_{\beta,r}^0$ ranks
below every QC of session $(\beta,r)$.
Recovery uses it when no QC of the session is known. A QC of another
session is never a justification in session $(\beta,r)$, so it cannot replace
$A_{\beta,r,0,0}$.

\paragraph{Regional receipts.}
After a replica commits a DATA block $B$ of session $(\beta,r)$, by $B$'s own
timer or as an ancestor, it sends
\[
 \langle\tagg{FINAL},\beta,r,\mathit{cfg}_r,h(B),\Hash(B)\rangle_v
\]
to its region. Matching FINAL messages of weight
$>F_{i_{\beta,r}}$ form the \defterm{regional receipt} $\FC_{\beta,r}(B)$. The receipt carries $B$ and the ancestry needed to validate it, or
refers to a validated local cache. Each correct holder of a complete
result offers it to the Publish Layer and multicasts the full validated
receipt, DATA and necessary dependencies to the fixed regional membership
on first acquisition. Deduplicate by session and DATA hash; retain the
outgoing action until dispatched. This rule also applies to a receipt
used for Completion and to a Completed or archival holder. It may also
send the result directly to other regions for early handoff.

\begin{definition}[\defterm{Valid regional receipt}]
\label{def:receipt}
A regional receipt $\FC_{\beta,r}(B)$ is valid when all of the following hold:
\begin{enumerate}[label=(\roman*)]
\item $B$ is a DATA block of session $(\beta,r)$, extends $A_{\beta,r,0,0}$, and has
$h(A_{\beta,r,0,0})<h(B)\le H_{\beta,r}$;
\item distinct members in session $(\beta,r)$'s fixed voting snapshot have
valid signatures on the same
$\langle\tagg{FINAL},\beta,r,\mathit{cfg}_r,h(B),\Hash(B)\rangle$,
with total weight $>F_{i_{\beta,r}}$ and the signed context specified above; and
\item $B$ and its required ancestry pass the block-validation rules of
\cref{sec:admission}.
\end{enumerate}
The verifier checks this evidence; it does not replay the signers'
timers. The weight threshold guarantees at least one correct signer,
who sends FINAL only after committing $B$.
\end{definition}
Here $\mathit{cfg}_r$ identifies session $(\beta,r)$'s \defterm{fixed voting snapshot}.
When we write $\FC(B)$ without a session subscript, we mean the receipt
from the session that produced $B$. The inherited anchor keeps that
receipt; it does not acquire a new one. A verifier outside
the region checks condition~(iii) from evidence, as
\cref{sec:admission} explains.
\begin{itemize}
\item FINAL binds no generation or term. So FINAL messages for one block
combine across recoveries. The block hash keeps the block's signed
context.
\item If a DRAIN commits $T_{\beta,r}$, replicas send FINAL for $T_{\beta,r}$, never for
the DRAIN.
\item Replicas of session $(\beta,r)$ send no FINAL for $A_{\beta,r,0,0}$. It keeps its original commitment evidence, which may come from an
earlier session if intermediate sessions added no DATA.
\item A replica may also commit $B$ and its uncommitted
ancestors when it obtains a valid $\FC_{\beta,r}(B)$ of its own session.
Completion does this for $T_{\beta,r}$ (\cref{alg:steady}). Such a commit
sends no FINAL, because the receipt already exists. Only Completion
also cancels timers. Receipts of other sessions serve only as history or anchor
evidence. They never change this session's locks or log.
\end{itemize}

\paragraph{Receipt coverage.}
\label{sec:coverage}
A valid $\FC_{\beta,r}(D)$ also proves commitment of every DATA ancestor $B$
of $D$ in the same session. At least one correct signer committed
$D$, which commits those ancestors. The \defterm{covering evidence} includes the
valid receipt for $D$ and the fully verified ancestry linking $D$ to $B$.
It does not claim that the signatures name $B$ directly.

This coverage never substitutes for a receipt from a different session.
It also does not replace the check that $B$ belongs to the authorized
external branch. A later Force may discard $D$ while keeping $B$; the
regional commitment evidence for $B$ still holds. External observers,
retained-anchor witnesses, and base justifications may use direct or
covering evidence. No extra FINAL messages or certificate round are
needed. We use ``\defterm{commitment evidence} for $B$'' for either form.

\subsection{Admission, validation, and engine states}
\label{sec:admission}
\paragraph{\defterm{Admission.}}
Replica $v$ admits views as follows.
\begin{itemize}
\item $(\beta,r,0,0)$ is admitted at $v$'s Ready event for the boundary that
opens session $(\beta,r)$ (\cref{def:ready}). The witness is the \defterm{entry proof}.
The NEWVIEW must choose exactly $A_{\beta,r,0,0}$.
\item $(\beta,r,g,0)$ with $g>0$ is admitted once $v$ holds
$\tagg{AUTH}_{r,g}$ and a quit proof for generation $g-1$. That proof is
the \defterm{entry proof}. The NEWVIEW is justified by a QC of an earlier
generation of session $(\beta,r)$, or by $Q_{\beta,r}^0$.
\item $(\beta,r,g,t)$ with $t>0$ is admitted once $v$ also holds the terminal
QCs of terms $0,\ldots,t-1$, each for DATA below $H_{\beta,r}$.
\end{itemize}
Term zero is admitted before its NEWVIEW arrives. Its progress checks
therefore run even if the initial leader is silent. Fresh proposals need
the generation's ANCHOR QC. This holds in every generation, including
generation $0$ of the first session from genesis; Part~I's
genesis shortcut does not apply. A
message is \defterm{admissible} if its view is admitted and its signatures
are valid.

A verifier that is not a replica of the session has no
admission events. When it validates history (\cref{def:receipt},
\cref{sec:external}), a view counts as admitted if the evidence contains
its \defterm{admission witness}: the boundary witness for $(\beta,r,0,0)$, the chain of
entry proofs for $(\beta,r,g,0)$, and the earlier terminal QCs for $t>0$.

\paragraph{\defterm{Complete validation.}}
Validation checks signatures, authority, admission, block type, quota,
both positions, the L1 and Espresso reference rules, the parent rule
(\cref{tab:parent}), parent QCs, hash links, and application validity
throughout the required ancestry. Every relevant ancestor includes its
signed proposal and block data; a chain of certificate hashes alone is
not enough. Dependencies come from the message or a validated local
cache. Correct senders send \defterm{complete validation bundles}. A reference is
allowed only if the receiver can resolve it within the delay bound of
\cref{ass:network}. Bundles list each dependency once.

\paragraph{Evidence takes priority.}
Before validating a message or discarding a duplicate, inspect all signed
material in it for a quit proof (\cref{sec:quit}). A proof needs its own
valid signatures and admission witnesses. A valid proof counts even if
the message that carries it is invalid or incomplete. Compare incoming
material with all retained material, not only with objects inside the
message.

\paragraph{One proposal-processing rule during normal operation.}
Unless retirement has closed the signing gate, every admitted fresh
proposal of the current generation that $v$ has not voted for invokes the
Vote step. Closing programs still inspect and forward signed evidence. This includes proposals in attached
histories, proposals of earlier terms, and proposals whose blocks are
already committed as ancestors. A proposal's identity is
$(\beta,r,g,t,j,\Hash(B))$. The signer set or encoding of its parent QC does
not change it. Deduplicate by this identity after evidence inspection.
Old-generation proposals are retained without new votes or timers.
Future-generation messages wait for entry.

\paragraph{Engine states.}
A replica runs one engine for each session pair $(\beta,r)$ assigned to its region.
Engines of different sessions remain separate, even if they use the
same session index. The engine of
session $(\beta,r)$ uses the following three core phases. The full
program lifecycle adds Waiting, Ready and Terminated in
\cref{sec:session-lifecycle}.
\begin{itemize}
\item \textbf{\defterm{Active}:} the current generation is admitted and unfrozen.
The engine proposes, votes, and changes terms.
\item \textbf{\defterm{Frozen}:} the current generation has a quit proof. The
engine stops proposals, votes, and term changes. It still processes
evidence, old votes, QCs, receipts, and recovery.
\item \textbf{\defterm{Completed}:} the replica holds $\FC_{\beta,r}(T_{\beta,r})$ satisfying
\cref{def:receipt}.
Production in session $(\beta,r)$ is over, even if the external handoff is still
pending. The engine only serves data.
\end{itemize}
Committing $T_{\beta,r}$ locally does not complete the engine. Other members may
still need this replica's recovery votes before a receipt forms. Each
message goes to the engine named by its signed pair $(\beta,r)$, so
evidence from one session or branch never freezes another. These engine
states do not describe external finality. The external chain may move
to another branch while an old engine still runs. L1 status
updates do not rewrite an engine's state or log. Effective Force closes
the external version; the manager then requests the explicit retirement
procedure of \cref{sec:retirement}. A closing Active or Frozen program
keeps that core phase until genuine Completion or FinishClose.

\begin{table}[htbp]\centering\small
\caption{Parent selection for the next fresh position $j$ of term $(\beta,r,g,t)$.}
\label{tab:parent}
\begin{tabularx}{\linewidth}{@{}l Y Y@{}}
\toprule
Condition & $B_{\mathrm{parent}}$ & $Q_{\mathrm{parent}}$\\
\midrule
$j=0,t=0$ (DATA or DRAIN) & Opening $A_{\beta,r,g,0}$ & Its ANCHOR QC in $(\beta,r,g,0)$\\
$j=0,t>0$ & Anchor $A_{\beta,r,g,t}$, the terminal block of $(\beta,r,g,t-1)$ & Its proposal QC in $(\beta,r,g,t-1)$\\
$1\le j<m_{\beta,r,g,t}$ & Position $j-1$ of $(\beta,r,g,t)$ & Its proposal QC in $(\beta,r,g,t)$\\
\bottomrule
\end{tabularx}
\end{table}

\section{Regional protocols}
\label{sec:protocols}
Protocols 1 and 2 run inside the engine of one session $(\beta,r)$. They are
Protocols 1 and 2 of Part~I~\cite{part1} with the changes listed below.
\RR{Red text} marks the retained regional changes from Part~I.
The new outer selection reference is defined in \cref{sec:blocks}.
External replay never changes the expiry of an existing regional timer.
Retirement blocks fresh signing, new timers and further entry under the
following gate; only genuine Quit or Completion cancels an existing timer.

\paragraph{Notation.}
The full view is $(\beta,r,g,t)$ throughout. Every anchor, receipt and
entry time below names its branch and session explicitly. In these two
protocols, $\Delta$ and $F$ mean $\Delta_{i_{\beta,r}}$ and
$F_{i_{\beta,r}}$. The authority fields $\mathit{cfg}_r$, $\TB_r$ and
$\tagg{AUTH}_{r,g}$ retain their numeric indices: a later branch that
reuses $r$ reads the same immutable authority. Its anchor, endpoint and
engine state can still differ.

\paragraph{Changes from Part~I.}
\begin{enumerate}[label=\arabic*.]
\item \emph{Stake.} Certificates need weight $>F$, not $f+1$ signatures.
\item \emph{Blocks.} A block also carries its session, physical
position, type, L1 reference $\rho$, and Espresso reference $\xi$
(\cref{sec:blocks}). Validity
checks all of them.
\item \emph{Quota.} A term has $m_{\beta,r,g,t}\le K$ positions. No term follows
$H_{\beta,r}$.
\item \emph{DRAIN.} A drain generation proposes one DRAIN. Its vote
counts as progress.
\item \emph{Receipts.} Replicas send FINAL after committing DATA. A
receipt for $T_{\beta,r}$ completes the engine (new step~5 of \cref{alg:steady}).
\item \emph{Timers.} Completion also cancels timers. A force never does.
\item \emph{Recovery entry.} After the $\Delta$ wait, a replica enters
when $\tagg{AUTH}_{r,g+1}$ is finalized. $Q_{\beta,r}^0$ can serve as the entry
lock and as the justification.
\item \emph{Genesis.} The initial session $(\beta,0)$ also certifies its anchor with an
ANCHOR QC before its first fresh proposal.
\item \emph{Sessions.} Each message goes to the engine of its session.
The initial session $(\beta,0)$ opens by bootstrap (\cref{alg:bootstrap}).
\end{enumerate}

\paragraph{Retirement gate and core events.}
Before any fresh proposal, DATA/DRAIN/ANCHOR vote, STATUS/NEWVIEW,
progress-check BLAME, term entry or recovery entry, require that
\texttt{close\_requested} is false and the program is not Terminated.
Apply the same check to buffered-message and authority-arrival callbacks.
This gate never prevents checking, comparing, retaining and relaying
signed evidence, processing a genuine Quit, executing a surviving commit
timer, sending its FINAL, or applying a valid receipt and Completion.
It does not overwrite a lock or cancel a timer. \Cref{sec:retirement}
defines its trigger and drain. Core entry, Quit, recovery, Completion
and timer-drain changes emit a \texttt{CoreEvent} to the local manager;
that event does not change public $G$.

\paragraph{Message-processing order.}\leavevmode
\RR{Replica $v$ first passes each message to the engine of the session
pair $(\beta,r)$ that the message binds. L1 updates go to \cref{alg:force}. Boundary witnesses and receipts of
other sessions go to \cref{alg:rotation,alg:external}.} The engine handles each received
message in this order:
\begin{enumerate}[label=(\arabic*)]
\item check all signed material for a quit proof;
\item check completeness and validity;
\item process attached DATA \RR{and DRAIN} proposals under Vote, parents
first;
\item process attached terminal QCs under step~4 of \cref{alg:steady},
in increasing term order;
\item process the message itself.
\end{enumerate}
When a message and a timer fall at the same time, $v$ handles the
message first.

\subsection{Steady state}
\label{sec:steady}
\Cref{alg:steady} gives the steady state of a session: proposals, votes,
timers, term rotation, and completion.

\begin{protocol}{alg:steady}{Propose, vote, commit, normal rotation\RR{, and completion}}
Replica $v$ runs these steps \RR{in the engine of session $(\beta,r)$}. Its
current term is $(\beta,r,g,t)$.
\begin{enumerate}[label=\textbf{\arabic*)},itemsep=3pt]
\item \textbf{Propose.} Run this step only if $v=L(\beta,r,g,t)$, $g$ is
unfrozen\RR{, and session $(\beta,r)$ is not complete}.
\RR{In a DATA generation,} for each position
$j=0,\ldots,\RR{m_{\beta,r,g,t}}-1$, in order:
 \begin{itemize}
 \item Wait until $v$ holds $B_{\mathrm{parent}}$ and
       $Q_{\mathrm{parent}}$ for $j$ (\cref{tab:parent}).
 \item Build \RR{$B=(\beta,r,k,h,\tagg{DATA},x,\Hash(B_{\mathrm{parent}}),\rho,\xi)$,
       with $k$ and $h$ from \cref{sec:quota} and the references
$\rho,\xi$ from \cref{sec:blocks}}.
       The batch $x$ follows \cref{sec:requests}.
 \item Broadcast
       $P=\langle\tagg{PROPOSE},\beta,r,g,t,j,B,Q_{\mathrm{parent}}\rangle_v$
       with its validation bundle \RR{to regional replicas. }
 \item Process $P$ under Vote.
 \end{itemize}
 \RR{In a drain generation, propose once, at $(t,j)=(0,0)$. The block is
 a DRAIN with parent $A_{\beta,r,g,0}$, no batch, DATA height $H_{\beta,r}$, and the L1
 and Espresso references of $T_{\beta,r}$. Broadcast and process it as above.}
 Do not wait for the parent to commit. Never sign two different
 blocks for one position. A resend repeats the same $P$.

\item \textbf{Vote.} Run this step on every DATA \RR{or DRAIN} proposal
$P$ of term $(\beta,r,g_P,t_P)$ at position $j_P$. Here $t_P$ may be earlier
than $t$. The step also applies to proposals inside attached histories.
 \begin{enumerate}[label=(\alph*)]
 \item If $P$ or its attached material contains a quit proof for $g$,
       run Quit and stop.
 \item If $g_P<g$, retain $P$ and stop. If $g_P>g$, buffer $P$
       until $v$ enters generation $g_P$.
 \item If $g$ is frozen \RR{or session $(\beta,r)$ is complete}, retain $P$ and
       stop.
 \item If a dependency is missing, retain $P$, request the dependency,
       and stop. Rerun this step when it arrives.
 \item If $P$ is invalid, discard it and stop. \RR{Validity includes the
       block type, the quota, both positions, and the L1 and Espresso references.}
 \item If $v$ has already voted for $P$, stop.
       A duplicate never restarts a timer.
 \item Otherwise, let $\tau_v(P)$ be the current time. Forward $P$ and
       its bundle to all \RR{regional replicas. }
       Broadcast
       $\langle\tagg{DATA\mbox{-}VOTE},\beta,r,g,t_P,j_P,\Hash(B)\rangle_v$
       \RR{(DRAIN-VOTE if $B$ is a DRAIN)}.
       Start the commit timer of $P$. It expires at $\tau_v(P)+2\Delta$.
 \end{enumerate}

\item \textbf{Commit.} When the timer of $P$ expires, commit $B$ and
its uncommitted ancestors. \RR{For each DATA block $D$ of session $(\beta,r)$
that this commits, send
$\langle\tagg{FINAL},\beta,r,\mathit{cfg}_r,h(D),\Hash(D)\rangle_v$ to the regional fixed membership and the Publish Layer collection interface. A committed DRAIN yields FINAL for $T_{\beta,r}$, never for
itself. } Only Quit \RR{and Completion} cancel a timer. Entering a
new term does not cancel it. A QC never starts a timer. \RR{A force
instruction never cancels a timer.}

\item \textbf{Terminal QC and rotation.} Run this step the first time
$v$ obtains a valid terminal QC for a term $(\beta,r,g,t')$.
``Obtains'' means that $v$ assembles the QC from votes or finds it in
any message, including as another proposal's $Q_{\mathrm{parent}}$.
 \begin{enumerate}[label=(\alph*)]
 \item Process each unseen proposal in the QC's history under Vote,
       parents first.
 \item If $g$ is now frozen, stop. \RR{If the QC certifies DATA at
       $H_{\beta,r}$, or a DRAIN, stop: no next term exists.}
 \item Send the QC and its validation bundle to $L(\beta,r,g,t'+1)$.
 \item If $t\le t'$, enter $(\beta,r,g,t'+1)$ and set $e_v(\beta,r,g,t'+1)$ to the
       current time. Cancel the progress checks of all earlier
       terms.
 \end{enumerate}
 A QC never moves $v$ to an earlier term. It never changes a recorded
 entry time.

\item \RR{\textbf{Completion.} Run this step the first time $v$ obtains a
regional receipt $\FC_{\beta,r}(T_{\beta,r})$ satisfying \cref{def:receipt}.}
 \begin{itemize}\color{rrRed}
 \item Multicast the complete validated terminal receipt and its dependencies
       to the fixed regional membership, once, and persist this outgoing action.
 \item Commit $T_{\beta,r}$ and its uncommitted ancestors.
 \item Mark session $(\beta,r)$ Completed. Stop its proposals, votes, progress
       checks, term changes, and recovery. Cancel its unexpired timers.
 \item Offer the complete terminal result to the Publish Layer.
       It may also be sent directly to the next region.
 \item Keep serving data. \Cref{alg:rotation} runs the handoff.
 \end{itemize}
 \RR{A force instruction never runs this step.}
\end{enumerate}
\end{protocol}

\Needspace{9\baselineskip}
The \defterm{acceptance time} $\tau_v(P)$ is when complete validation and the first
vote occur, not an earlier arrival of an incomplete packet. Each replica
records its own entry time for each term. In session $(\beta,r)$, $e_v(\beta,r,0,0)$ is
set at the atomic Ready/Enter event, in step~1 of \cref{alg:bootstrap}. For $g>0$, $e_v(\beta,r,g,0)$ is
set in step~3 of \cref{alg:recovery}. For $t>0$, $e_v(\beta,r,g,t)$ is set in
step~4 of \cref{alg:steady}. A term's progress checks therefore start
when the term is admitted locally, not when its leader's first proposal
arrives. A silent new leader cannot delay them.

A proposal or QC at height $H_{\beta,r}$ does not end the session. Without a
quit, the timer of $T_{\beta,r}$ commits it; no later block is needed. Only a
receipt for $T_{\beta,r}$ completes the engine.

\subsection{Quit proofs and recovery}
\label{sec:quit}
A quit proof freezes a generation. \Cref{alg:recovery} then moves the
session to the next generation.

\begin{protocol}{alg:recovery}{Blame, quit, and recovery}
Replica $v$ runs these steps \RR{in the engine of session $(\beta,r)$}. Its
current term is $(\beta,r,g,t)$.
\begin{enumerate}[label=\textbf{\arabic*)},itemsep=3pt]
\item \textbf{Blame.} \RR{Let $q=m_{\beta,r,g,t}$ in a DATA generation and
$q=1$ in a drain generation.} For each $p=1,\ldots,\RR{q}$, run a
progress check at time
\[
 e_v(\beta,r,g,t)+(2p+4)\Delta.
\]
\RR{Schedule the check for $p=1$ at entry, even if the opening is not yet
known. Schedule the others once $q$ is known. A check whose time has
already passed runs at once.}
At that time, broadcast $\langle\tagg{BLAME},\beta,r,g,t\rangle_v$ if all of
these hold:
 \begin{itemize}
 \item $v$'s current term is still $(\beta,r,g,t)$, $g$ is unfrozen\RR{, and
       session $(\beta,r)$ is not complete};
 \item $v$ has cast DATA \RR{or DRAIN} votes for fewer than $p$ positions
       of term $(\beta,r,g,t)$ (ANCHOR votes do not count);
 \item $v$ has not yet sent BLAME for $(\beta,r,g,t)$.
 \end{itemize}

\item \textbf{Quit.} Run this step the first time $v$ holds a quit proof
for its current generation $g$.
 \begin{itemize}
 \item Freeze $g$. Stop all progress checks, proposals, votes, and
       term changes in $g$. Cancel every unexpired commit timer of $g$.
 \item Set $a_v$ to the current time. Broadcast the proof.
 \item Keep all commitments. Keep collecting votes and QCs of $g$.
       They can raise the entry lock chosen in step~3. \RR{Keep
       processing regional receipts (step~5 of \cref{alg:steady}).}
 \item If the proof is of type (b) or (c), submit it to L1. Do
       not wait for L1 to process it.
 \end{itemize}

\item \textbf{Status and entry.} \RR{Skip this step if session $(\beta,r)$ is
complete.} \RR{Wait until at least $a_v+\Delta$ and until
$\tagg{AUTH}_{r,g+1}$ is finalized. Set $e_v(\beta,r,g+1,0)$ to the time both
conditions hold, and read that authority from L1.} Then:
 \begin{itemize}
 \item Set the entry lock $Q_v$ to the highest-ranked QC
       (DATA\RR{, DRAIN,} or ANCHOR) that $v$ knows from generations
       $0,\ldots,g$\RR{, or to $Q_{\beta,r}^0$ if it knows none}.
       $Q_v$ never changes after this. QCs that arrive later do not
       update it.
 \item Enter $(\beta,r,g+1,0)$ as unfrozen and start its progress checks.
       \RR{Keep $A_{\beta,r,0,0}$, $H_{\beta,r}$, and the voting snapshot of session $(\beta,r)$.}
 \item Send $\langle\tagg{STATUS},\beta,r,g+1,0,Q_v,E\rangle_v$ and its
       validation bundle to $L(\beta,r,g+1,0)$. Here $E$ is the quit proof for
       $g$, which is also the entry proof of $g+1$.
 \item Process any buffered proofs for $g+1$.
 \end{itemize}

\item \textbf{New-view.} Run this step only at $L(\beta,r,g+1,0)$. Wait $2\Delta$
after its own entry. If $g+1$ is still current and unfrozen:
 \begin{itemize}
 \item Let $Q^*$ be the highest-ranked QC from generations $0,\ldots,g$
       that the leader knows, including the QCs in STATUS
       messages\RR{; if it knows none, let $Q^*=Q_{\beta,r}^0$}.
 \item Let $A_{\beta,r,g+1,0}$ be the block that $Q^*$ certifies. \RR{If
       $h(A_{\beta,r,g+1,0})<H_{\beta,r}$, generation $g+1$ is a DATA generation.
       Otherwise it is a drain generation.}
 \item Broadcast
       $\langle\tagg{NEWVIEW},\beta,r,g+1,0,A_{\beta,r,g+1,0},Q^*,E\rangle_{L(\beta,r,g+1,0)}$
       with its validation bundle. Choose only one opening block for
       $g+1$.
 \end{itemize}

\item \textbf{Anchor vote.} On receiving a NEWVIEW for $g+1$:
 \begin{enumerate}[label=(\alph*)]
 \item If it contains a quit proof for $g+1$, run Quit and stop.
 \item Validate its leader, entry proof $E$, opening block $A_{\beta,r,g+1,0}$,
       $Q^*$, and dependencies. If data are missing, request them and
       wait without voting. If the message is invalid, discard it.
 \item Vote only if all of these hold: $g+1$ is current and unfrozen; $v$
       has cast no ANCHOR vote in $g+1$; and $\rank(Q^*)\ge\rank(Q_v)$.
 \item To vote, forward the NEWVIEW to all replicas. Broadcast
       \[\langle\tagg{ANCHOR\mbox{-}VOTE},\beta,r,g+1,0,\Hash(A_{\beta,r,g+1,0})\rangle_v.\]
       Start no timer. Use no fresh position.
 \end{enumerate}
 Matching ANCHOR votes \RR{of total weight $>F$} form the ANCHOR QC.
 When $L(\beta,r,g+1,0)$ obtains it, the leader starts Propose in
 \cref{alg:steady} at $j=0$\RR{; in a drain generation, this is the
 DRAIN}. Every replica, including one that cast no ANCHOR vote,
 processes DATA \RR{and DRAIN} proposals backed by a valid ANCHOR QC
 under Vote.
\end{enumerate}
\end{protocol}

\paragraph{\defterm{Complete quit proofs.}}
For an admitted generation $g$ \RR{of session $(\beta,r)$}, a proof consists of
either:
\begin{enumerate}[label=(\alph*)]
\item distinct signatures \RR{of total weight $>F$} on the same
$\langle\tagg{BLAME},\beta,r,g,t\rangle$ for an admitted term, with that term's
admission witness;
\item two different admissible DATA \RR{or DRAIN} blocks signed by
$L(\beta,r,g,t)$ at the same $(\beta,r,g,t,j)$, with their admission \RR{and quota}
witnesses; or
\item two valid NEWVIEWs signed by $L(\beta,r,g,0)$ selecting different opening
blocks for generation $g$.
\end{enumerate}
\RR{A DATA block beyond the quota, or a DRAIN outside a drain generation,
is invalid. It is rejected and never counts for (b).}
Different openings count even when their blocks are comparable, because
they fix different starting heights $h(A_{\beta,r,g,0})$. The height formula then
maps the same position to different heights.

Different QCs for the same opening block are not equivocation. Evidence
from any admitted term can freeze its generation, including a term not
yet locally reached. Old-generation proofs are archival and may be
submitted to L1. \RR{Evidence from another session never freezes this
engine.} Future-generation proofs are buffered until entry. Neither
resets an unrelated recovery.

The Publish Layer uses these DATA and FINAL receipts. No publication
acknowledgement is required for regional commitment or Completion.
Participants may forward regional bundles directly; no external send
or Espresso confirmation delays a regional vote or timer.

The entry-lock test governs ANCHOR voting only. A replica that has
committed $T_{\beta,r}$ locally still takes part in recovery until it holds
$\FC_{\beta,r}(T_{\beta,r})$, because other members may need its votes. This continued-participation
rule applies until a verified RequestClose closes the fresh-action gate.

\section{Regional and global finality}
\label{sec:external}
Regional commitment and global acceptance are separate events.
A node follows a production branch $R_u$ and a continuous global DATA
prefix $G_u$. The branch may extend several normal sessions beyond
$G_u$. A Force closes that extension, never $G_u$ itself.

\subsection{Regional finality}
\label{sec:regional-finality}
The commit-window progress statements below concern normally running
sessions with no RequestClose during the stated interval. Closing
sessions instead use \cref{sec:retirement}; they may retain a local
commitment without forming a complete receipt. This qualification is a
scope of the claim, not an added assumption about all executions.
\begin{definition}[\defterm{Regional finality}]
\label{def:regional-finality}
A DATA block $B$ has regional finality on a specified production branch
when its boundary history assigns it to its claimed region and session,
it has valid direct or same-session covering commitment evidence, and
the DATA path to it is valid and continuous. Each session on that path
supplies its own authority and commitment evidence.
\end{definition}
The region is derived from the boundary before its receipt is checked.
A valid receipt does not authorize another region's DATA or a different
session. Current live/off state does not replace historical authority.
When a Force closes an unconfirmed suffix, the receipts remain evidence
about the old regional engine. They no longer authorize that suffix to
extend the global prefix.

\paragraph{Remark (Commitment by one replica and by all replicas).}
Here ``all replicas'' means all correct voters of the session that
produced $B$. Consider that session's regional engine, with recovery
authorities set in advance.
\begin{enumerate}[label=(\arabic*),itemsep=4pt]
\item \textbf{What a receipt proves.} A valid direct or covering
receipt proves that at least one correct replica has committed $B$.
Other correct replicas may still have running timers, or may have
canceled their timers on Quit. Regional finality does not assert that
they have all committed already.

\item \textbf{When all replicas commit.} Regional safety preserves
$B$ through recovery: the selected anchor is $B$ or a descendant of
$B$. ANCHOR voting re-certifies that block but starts no commit timer.
Suppose a correct replica votes for and forwards a fresh DATA or DRAIN
proposal for a descendant $D$ at time $s$. If no correct replica freezes
that generation through $s+3\Delta_{i_{\beta,r}}$, all correct replicas commit
$B$ by that time. Replicas that have not completed the session vote by
$s+\Delta_{i_{\beta,r}}$; their timers commit $D$ and its ancestors at most
$2\Delta_{i_{\beta,r}}$ later. Completion may cancel a timer, but first commits
$T_{\beta,r}$ and its ancestors, including $B$. In particular, a DRAIN need not
itself commit at a replica that completes through the receipt for $T_{\beta,r}$.

\item \textbf{Why one honest leader is not enough.} An honest leader
may finish its term before its timers expire. A Byzantine successor
can then equivocate and freeze the generation. What is needed is the
uninterrupted commit window in (2). Under regional safety and the
progress conditions of \cref{sec:progress}, including permanent
equivocation exclusion and fair leader schedules, the engine eventually
obtains such a window while it continues to run, or completes through
the terminal receipt. Earlier recoveries do not revoke $B$. These
conditions do not give a fixed bound on when the window begins.

\item \textbf{When a replica can know.} A correct regional replica
that directly commits $D$ by its timer at time $c$ can conclude that
all correct regional replicas have committed $B$ if it remains Active
in the same generation through $c+2\Delta_{i_{\beta,r}}$, after processing
messages due at that time. Other replicas' timers are due by
$c+\Delta_{i_{\beta,r}}$. A Quit that canceled one of those timers would reach
this replica by $c+2\Delta_{i_{\beta,r}}$ and freeze its generation. A replica
whose timer was canceled by Completion has already committed $B$.
Normal term rotation does not affect this test.

\item \textbf{What an external observer can know.} A receipt alone
does not establish that uninterrupted window. The delivery bounds cover
complete submitted results, not every regional BLAME or quit event. The
absence of a published quit proof therefore does not prove that no replica
quit. Espresso evidence establishes public availability of the result.
A bound on when every replica executes its local commit step would also
need the relevant online and receipt-processing conditions. This protocol
uses Espresso acceptance to protect the external prefix against Force;
it does not use an external all-replicas-committed test.
\end{enumerate}



\subsection{Published records and the global interpreter}
\label{sec:exports}
The Publish Layer of \cref{sec:delivery} carries complete regional
results to Espresso. A sender is an existing participant performing this
task, not a new consensus authority. It builds an EXPORT, submits it and
tracks whether its DATA enters $G$. The leader-first and sponsored-backup
implementation is deferred to \cref{sec:publish-implementation}.

An \tagg{EXPORT} record binds the chain identifier, branch version,
expected global base, an ordered nonempty interval of DATA, and the
evidence needed to validate that interval. Evidence includes regional
receipts, ancestry, REGULAR boundaries, historical authorities and
native application context. A direct or covering receipt is supplied
for each session's DATA. A receipt from a later session does not cover
an earlier session. DRAIN may occur in supporting regional ancestry,
but is never exported as application DATA.

An EXPORT can contain the missing end of region A and the beginning
of region B. It therefore need not wait for a separate L1 publication
or for each ancestor to be submitted in its own Espresso transaction.
The entire record is checked on a temporary state. Either all of its
DATA, counts and normal boundary transitions are accepted, or none are.
Another receipt encoding or signer set does not create another DATA
position. Publication and global acceptance are different events:
$T_{\rm pub}+T_E$ bounds delivery of the submitted result, not acceptance
of an invalid or stale EXPORT. An authenticated record can still provide
regional evidence for direct handoff or replay even if that particular
EXPORT does not advance $G$. Its contents must pass their own checks.

\paragraph{Acceptance.}
After the block's prelude in \cref{alg:force}, a record is accepted
only if all of these checks succeed:
\begin{enumerate}[label=(\roman*)]
\item the expected base equals the current $G$, including its prefix
commitment and session provenance;
\item the record's branch version is open, and every block is assigned to the
session reached by the validated normal boundary chain;
\item the DATA heights begin at $h(G)+1$, parents and hashes agree,
and no input or intermediate session is omitted or repeated;
\item regional commitment evidence, application validity and quota
checks succeed for every block;
\item each Espresso selection reference used by the record lies
strictly before its containing Espresso block, and its historical
eligibility snapshot is authenticated;
\item a branch version created by a Force in $J_q$ first exports in a block
strictly after $J_q$.
\end{enumerate}
The record is then committed atomically in namespace order.
Each new DATA increases the appropriate export-session count. On its
terminal DATA, that session's export window is cancelled. Its normal
successor becomes the next export session when the fixed selection
rule resolves. Its regional engine may have started much earlier.

A valid terminal DATA may enter $G$ even if no successor is yet
eligible. The interpreter keeps that terminal anchor and enters
\defterm{unassigned}. It runs no DATA window while no export duty is
assigned. A proven absence of a candidate is different from unavailable
evidence: missing history or authority evidence pauses the relevant
verification and is never treated as an empty candidate set.
When deferred evidence is obtained, replay resumes at the same canonical
position; local retrieval time never supplies a new public timestamp.

Malformed records and well-formed but invalid candidates have no
effect on $G$ or its counts. A future-height record skipped before its
ancestors are accepted must be submitted again. The interpreter never
goes back to make an earlier rejected record global. An incomplete
namespace download is not a malformed record: it prevents the block
from being processed at all.

\subsection{Global finality}
\label{sec:global-finality}
Let $S_u$ be observer $u$'s verified interpreter state. It contains $G_u$,
the current branch version and export duty, windows, eligibility state, pending
controls, and source-verification checkpoints.
Initially it is a trusted checkpoint or the state derived from genesis.

\begin{definition}[\defterm{Global finality at an observer}]
\label{def:global-finality}
Observer $u$ treats DATA block $B$ as globally final when it has
authenticated and replayed a finalized Espresso prefix containing an
accepted EXPORT that adds $B$ to $G_u$. Replay uses complete applicable
L1 control history and the event order of \cref{alg:force}. All DATA
predecessors are already accepted or are accepted in the same EXPORT.
\end{definition}

No extra waiting period follows this verification. A later Force
preserves the entire existing $G_u$, including its original block and
session identities. An old-version record after closure can still be
finalized as raw Espresso bytes, but it is not an accepted chain result.
Espresso inclusion alone, or a receipt alone, does not establish this
definition.

This is input-prefix finality. Espresso does not execute the native
transactions or prove their resulting state. The regional validity
contract and the execution-verification interface remain responsible
for that property. Finality of a state claim or withdrawal on L1 is a
separate stage. Publication lag never permits skipping globally final
input.

\subsection{References do not provide a commit clock}
\label{sec:late-receipt}
A block's L1 reference and Espresso selection reference are fixed when
it is proposed. A later recovery or DRAIN may commit it long afterwards.
Neither old reference starts an export window or establishes global
finality. An export window starts when replay assigns the next DATA
duty, and counts only newly accepted DATA from finalized Espresso
records. Regional timers continue to use the real-time clock of
\cref{ass:network}, not Espresso timestamps.

\section{Cross-region rotation}
\label{sec:xprotocols}

\subsection{Four-module cross-region architecture}
\label{sec:runtime-modules}
Each correct regional replica runs a \defterm{Session Manager}. The
manager observes finalized L1 and Espresso history, receives regional
results, and manages independent programs named by $(\beta,r)$.
It is local replicated logic, not an additional trusted regional
controller. A region can run an old session while preparing a possible
successor and handing an earlier session's data duties to storage and
publication workers. The full regional view remains $(\beta,r,g,t)$.

The cross-region interface has four modules. The Soft Force Detector
recomputes delivery failure from authenticated history. The Multisig
Health Controller submits the existing L1 controls from health and
operational observations. These are parallel sources of intervention.
The Cross-region Handler manages candidate and fixed handoffs. Each
session's Ready and Activation Gate checks its exact authorized starting
point before entering the regional engine.

\begin{table}[htbp]\centering\small
\caption{The four cross-region modules.}
\label{tab:runtime-modules}
\begin{tabularx}{\linewidth}{@{}>{\raggedright\arraybackslash}p{.22\linewidth}Y Y@{}}\toprule
Module and executor & Inputs & Output and authority limit\\\midrule
Soft detector; any verifier & Complete verified Espresso history,
export duty, fixed windows and accepted counts & Recomputable failed-window
condition; it does not itself change $\beta$.\\
Multisig health controller; the existing committee & Regional health,
synchronization and delivery observations & Future L1 MSIG, STATUS or
CLEAR; it does not choose a retained height.\\
Cross-region handler; each correct regional replica & Verified public
updates, regional receipts and local task events & Prepare, activate,
request close and transfer responsibilities; it does not replay Force
a second time.\\
Ready and activation gate; each session program & Fixed handoff,
exact anchor, execution state and authority & Waiting or atomic entry
through \cref{alg:bootstrap}; prefetching grants no voting right.\\\bottomrule
\end{tabularx}
\end{table}

Healthy regions keep watching both public streams even while they have
no Active session. They can prepare data and authorities before being
selected. Their public replay cursor, provisional programs and storage
tasks are separate objects. Local lifecycle phases never determine
whether an EXPORT may advance $G$.


\begin{figure}[p]\centering
\includegraphics[width=.98\linewidth,height=.87\textheight,keepaspectratio]{figures/cross_region_rotation_v3_2/modules.pdf}
\caption{The four cross-region modules and their shared verification and
publication services. Each regional node manages independent programs;
prewarming and fixed-handoff activation use distinct paths.}
\label{fig:runtime-modules}
\end{figure}

\subsection{Soft detection and multisig health control}
\label{sec:runtime-controls}
The Soft Force Detector takes a validated replay checkpoint and the
complete following Espresso prefix. It reproduces the current export
duty and settles its fixed windows as in \cref{sec:windows}. It returns
either no failure, a failed-window descriptor with its first crossing
header and supporting history, or a need for missing evidence. The
descriptor binds the chain, current branch version, export-session pair,
window start, deadline, target and accepted count. Missing namespace or
control evidence is never a zero count.

Only DATA actually appended to $G$ counts. Duplicates, DRAIN, proposal
QCs and invalid EXPORTs do not count. The observed export session can
lag several speculative production sessions. Polling only supplies new
history; its frequency and local reception time do not decide a failure.
Every recipient recomputes a supplied descriptor against its own
authenticated prefix. Protocol~5 remains the sole Force transition:
independent STATUS/CLEAR, the first matching MSIG, otherwise the
automatic check, then the block's EXPORT payload. A computed soft
condition cannot bypass a matching hard intervention in that prelude.

The committee's health policy distinguishes four actions. To exclude
an unsuitable region from future selection, it submits a future OFF.
To end the current target version's production arrangement, it submits
a future MSIG Force, with an OFF if continued exclusion is intended.
After checking synchronization, data service and operational recovery,
it may submit ON for future selections. To remove an automatic cooldown
before expiry, it submits CLEAR explicitly. ON alone does not clear it.
Each operation records its target, authorized changes and prospective
activation time; concrete health thresholds remain an operational
choice under \cref{ass:committee}.

An MSIG targets a branch version. A regular handoff changes $r$ and
therefore does not invalidate it. An earlier Force changes $\beta$ and
can make it stale. Its attached configuration applies only when the
target matches; an independent STATUS follows its own activation rule.
Health observations do not directly change another replica's public
state. All actions pass through finalized L1 and their future Espresso
activation positions under the existing announcement condition.
An ON does not certify future regional synchrony or replace the
regional fault assumptions.

\subsection{Cross-region Handler}
\label{sec:session-manager}
The manager serializes events that modify its local control state.
Data retrieval and state restoration may run concurrently. Its
inputs and outputs have the following meanings.

\begin{table}[htbp]\centering\small
\caption{Session Manager inputs.}
\label{tab:manager-inputs}
\begin{tabularx}{\linewidth}{@{}>{\raggedright\arraybackslash}p{.24\linewidth}Y@{}}\toprule
Event & Validation and routing\\\midrule
\texttt{L1Finalized} & Authenticate the complete control prefix and
authorities; queue future controls without activating them.\\
\texttt{EspressoFinalized} & Authenticate complete ordered history;
consume the state transition produced by \cref{alg:external,alg:force}.\\
\texttt{RegionalFinalized} & Validate the DATA, complete FINAL receipt
and required ancestry; route to the producing pair and offer the result
to the Publish Layer.\\
\texttt{MaterialAvailable} & Check the target pair, anchor's producing
pair, height and hash, requested object and verified result; update only
the matching preparation or fixed handoff.\\
\texttt{CoreEvent} & Accept an event from the named local engine;
record its view, quit, recovery, Completion or retirement progress.\\\bottomrule
\end{tabularx}
\end{table}

One FINAL signature is not a \texttt{RegionalFinalized} event. Existing
participants collect matching signatures into the receipt of
\cref{def:receipt}. A correct holder delivers the complete result to
the manager and the entered source engine, multicasts the complete
receipt and dependencies to the fixed regional membership, and offers
it to the Publish Layer. The source engine applies the existing receipt and Completion
rules. A program that has not entered only caches the result; receipt
arrival cannot create an engine or skip its handoff checks. Aggregation
is not a new consensus round. $T_{\rm pub}$ still starts when a complete
result becomes available.

The outputs are \texttt{Prepare/Fetch}, naming the target and exact
requested material; \texttt{Activate}, invoking Protocol~4 once;
\texttt{RequestClose}, carrying a verified closure cause;
\texttt{OfferResult}, preserving original DATA and receipt fields;
and \texttt{Persist/Release}, recording state and transferring resource
responsibilities. A cached callback is material, not an activation
command. A late result can be useful to storage without matching the
program's current preparation target.

\begin{protocol}[H]{alg:rotation}{Handle cross-region events and manage sessions}
Each correct regional replica runs this handler independently.
\begin{enumerate}[label=\textbf{\arabic*)},itemsep=3pt]
\item \textbf{Validate the event.} Route authenticated public input,
regional results and local callbacks by their stated context. Preserve
pending requests and source evidence. Reject a purported completion
whose material or session does not match its request.
\item \textbf{Consume public replay.} For new public history, run or
consume \cref{alg:external}, including \cref{alg:force} in its canonical
order. Use its resulting $\beta$, $G$, eligibility snapshots, duty and
boundary records. This handler never increments $\beta$, chooses a
different cutoff or modifies $G$ independently.
\item \textbf{Reconcile handoffs.} From valid terminal receipts,
derive regular boundaries while their version remains open. Consume
fixed Force boundaries from replay. Verify them by
\cref{sec:witness}. Prepare selected local targets; send complete
handoff material to the selected members when available. A future
Force does not block a current regular handoff.
\item \textbf{Reconcile candidates.} Recompute possible future
targets from the current export context, $G$, eligible regions and
known future controls. Update only unfixed preparations. A different
target pair uses a different record. Pause obsolete speculative work;
terminate an unentered program only when verified facts make that
target impossible or assign its fixed handoff to another region.
\item \textbf{Handle closure and results.} For an entered program
whose version is closed, issue the idempotent \texttt{RequestClose}
of \cref{sec:retirement}. Keep its signed routing, historical evidence
and publication duties. Apply terminal receipts to their entered
source engines under the original Completion rule.
\item \textbf{Check and enter.} For each selected local Waiting
program with a fixed handoff, check \cref{def:ready}. In the same local
event, recheck public validity and the durable entry flag, record Ready,
and invoke Protocol~4 Enter atomically. A failed check leaves Waiting
or applies a justified cancellation; it does not partially initialize
the engine.
\item \textbf{Persist and dispatch.} Durably record state changes,
entry and exit records, and required outgoing actions before exposing
them. Dispatch asynchronous work with its full target and anchor
context. Repeated events resume outstanding actions without creating
another logical entry.
\end{enumerate}
\end{protocol}

The public interpreter has one mutation path: Protocol~5 defines
control, window and Force transitions; Protocol~6 authenticates and
applies them and accepts EXPORTs. Protocol~3 consumes those results.
Protocols~1 and~2 update only the named session's regional core.
An implementation may place these modules in one process without
changing this division of authority.

\subsection{Session lifecycle}
\label{sec:session-lifecycle}
The manager maintains the following record for each local program:
\[
\begin{aligned}
\tagg{Program}[\beta,r]=(&\mathit{phase},\mathit{preparing\_plan},
 \mathit{fixed\_handoff},\mathit{anchor\_block},\\
 &\mathit{execution\_state},\mathit{authority\_and\_bounds},
 \mathit{core\_state},\mathit{pending\_fetches},\\
 &\mathit{entered},\mathit{close\_requested},\mathit{closure\_evidence},\\
 &c_v,d_v,\mathit{publication\_and\_retention\_handles}).
\end{aligned}
\]
The six phases describe the local program, not public finality.
\begin{description}[style=nextline,leftmargin=1.5em,labelindent=0pt,labelwidth=1.5em]
\item[Waiting.] The program is preparing a possible or selected
handoff. It may fetch data, recover state and cache evidence. It has
not entered its regional engine, signed a regional vote or started
a regional commit or progress timer.
\item[Ready.] A fixed handoff, its exact anchor state and all entry
dependencies have passed the gate below. Enter follows in the same
serialized event; Ready introduces no extra scheduling delay.
\item[Active.] The program has entered an unfrozen regional
generation. Initial STATUS, the $2\Delta$ wait, NEWVIEW and ANCHOR
certification all occur in Active. The core runs Protocols~1 and~2,
subject only to the explicit retirement interface in
\cref{sec:retirement} when closure is requested.
\item[Frozen.] The current generation has a valid Quit proof.
The original evidence, receipt and recovery rules apply. Force alone
does not establish a Quit proof or cause this phase.
\item[Completed.] The entered engine has obtained a valid terminal
receipt and performed the original Completion action. It has stopped
regional production and cancelled timers as specified there. Data
delivery and retention responsibilities may remain.
\item[Terminated.] The logical program has exited. It processes no
BFT timer and issues no BFT signature. Durable storage and publication
workers retain its required evidence and duties. An exit record prevents
late messages from restarting it.
\end{description}

Normal transitions are Waiting to Ready to Active; Active to Frozen
on Quit; Frozen to Active on the original recovery admission, with
$g$ advanced and $t=0$; and Active or Frozen to Completed on terminal
Completion. Waiting or Ready may terminate when their target is
definitively invalid. Completed may terminate after responsibility
transfer; if RequestClose began while Active or Frozen, the original
drain deadline still applies. Active or Frozen may terminate only through the retirement
contract in \cref{sec:retirement}. While close is pending, they retain
their core phase and set \texttt{close\_requested}; there is no seventh
phase.

An unfixed \texttt{preparing\_plan} can change when $G$, the source
duty or candidate eligibility changes. The fixed handoff cannot.
Its target pair, assigned region, initial anchor, source context,
selection position, authority and reference bounds remain immutable.
In particular, a later OFF does not replace a historical regular
selection. Local locks, votes, view entry times and timers belong to
one pair only. Recovering a crashed process resumes this durable
logical record; it does not clear its votes or perform another initial
Enter. Storage and task references carry the full pair and anchor
identity. No additional protocol-level session identifier is needed.

Several regions may hold unfixed Waiting preparations for the same
future pair. Only the region selected by the fixed handoff can enter.
Temporary loss of candidate eligibility may pause work without
terminating a preparation that could become eligible again. A retired
record remains a tombstone even after active resources are released.

An Active old program may already lack permission to extend $G$.
A Completed or Terminated program's existing results may still be
accepted if their version and interval satisfy the public acceptance
rules. The phase does not revoke or create that permission.

\subsection{Handoff records, anchors and Ready}
\label{sec:witness}
Let $\tagg{NextLive}(i,M)$ first search the other regions after $i$
in the fixed cyclic order $\mathcal L_R$. If none is eligible, return
$i$ itself when it is eligible; otherwise return $\bot$. Eligibility
still combines the historical governance flag, cooldown and positive
designated-authority weight. For the synthetic genesis source $\bot$,
search the whole list from its first region. Same-region renewal creates
a new pair and fresh core state; it never reuses the old engine.

An AUTO cooldown still applies to its source. If no alternative exists,
the pending boundary waits until the source becomes eligible again;
duplicate replay never restarts that cooldown. OFF remains effective
until a later ON. An MSIG that intends to exclude its source must include
or separately schedule OFF; Force alone may renew that same region when
it is the only eligible candidate. This is an explicit selector change,
not a new availability premise.

A pending boundary stores source context
$(i_b,\beta_b,r_b)$ and next index $r_b+1$. A regular boundary uses
its completed predecessor. A Force with an assigned export duty
captures that duty's context. With no assigned duty, a Force reuses
the stored source context. Repeated Forces before assignment replace
the pending cause and version without incrementing an unopened
session. Genesis uses $(\bot,\bot,-1)$.

For REGULAR, the successor of $(\beta,r)$ is $(\beta,r+1)$ and
inherits exactly $T_{\beta,r}$. For Force, the new version's index is
$r_b+1$ and its anchor is exactly $G$ before the triggering block's
payload. The producer of that anchor need not be the captured export
session. The anchor keeps its original producing pair, hash and
commitment evidence.

For a regular boundary set $q_0$ to the height of
$\xi(T_{\beta,r})$; for Force set $q_0$ to its triggering height.
Using the source region just defined and the designated next authority,
the fixed selection is
\[
 q_s=\min\{q\ge q_0:
 \tagg{NextLive}(i_b,\chi_E(q))\ne\bot\},\qquad
 i_{\rm next}=\tagg{NextLive}(i_b,\chi_E(q_s)).
\]
The search uses verified history only. Missing evidence causes a wait,
not a skipped candidate. A later snapshot cannot change an established
minimum. Regular selection need not wait for the terminal DATA to be
published. Its signed selection reference may be historical; no
receive-time freshness rule is added here.

\paragraph{Typed handoff packet.}
A handoff packet contains the fields below. References can name
authenticated cached objects or objects still to be fetched; a reference
alone does not establish Ready.
\begin{table}[htbp]\centering\small
\caption{Handoff packet fields.}
\label{tab:handoff-packet}
\begin{tabularx}{\linewidth}{@{}>{\raggedright\arraybackslash}p{.26\linewidth}Y@{}}\toprule
Field & Type and meaning\\\midrule
Domain and kind & Chain/protocol domain; one of REGULAR, MSIG, AUTO.\\
Source context & Predecessor pair $(\beta_b,r_b)$ and source region
$i_b$, including the stored context when no duty is assigned.\\
Target & Target pair $(\beta',r')$ and selected region $i_{\rm next}$.\\
Anchor & Producing pair, DATA height, hash, full DATA or reference,
and original direct or covering commitment evidence.\\
Authority & Finalized $\tagg{AUTH}_{r',0}$, its fixed configuration,
quota, and publication-height proof.\\
Selection and bounds & $q_0,q_s$, authenticated selection snapshot,
$q^{\rm open}_{\beta',r'}$ and $\lambda^{\rm open}_{\beta',r'}$.\\
Cause and evidence & Mode-specific cause; trusted or validated replay
checkpoint, complete required L1/Espresso history and regional proofs.\\\bottomrule
\end{tabularx}
\end{table}

\paragraph{Handoff verification.}
\label{sec:verify-handoff}
The pure procedure $\tagg{VerifyHandoffHistory}(C,\mathcal H,P)$ takes an
authorized replay checkpoint $C$, required source history $\mathcal H$
and packet $P$. It returns historical fixed fields, invalid, or needs
evidence. It authenticates what the boundary established; it does not
require its version to remain open now and grants no entry permission.
A separate activation predicate is given after the procedure. Neither
procedure changes public replay.
\begin{enumerate}[label=(\roman*),itemsep=2pt]
\item Check the domain, field types, source and target names. Replay
authenticated compatible history from $C$ through all positions
needed by the cause, historical selection and original admission.
Verify complete applicable controls, not only inclusion of the named
control. Consult validated caches when available.
\item For REGULAR, validate the predecessor's own fixed boundary,
initial authority, exact quota and terminal DATA
$h(T_{\beta_b,r_b})=h(A_{\beta_b,r_b,0,0})+\TB_{r_b}$.
Validate its terminal FINAL receipt, ancestry and signed $\xi$.
Require $\beta'=\beta_b$ and $r'=r_b+1$. The anchor must equal the
terminal DATA. Recursively validate the predecessor boundary in this same
historical mode; current closure does not erase its original evidence.
\item For MSIG, verify its authorized finalized L1 record, control
sequence and first-crossing header $J_{q(\theta_a)}$. Replay the
entire prelude and check that this is the first due Force whose target
matches the entry version. Check attached changes only under this
match. A stale instruction or another due instruction consumed by the
same prelude cannot form a second boundary.
\item For AUTO, replay the fixed window's count and first deadline
crossing from complete history. Check insufficient accepted DATA and
the absence of a matching higher-priority MSIG in that prelude.
Check the cooldown derived at the realized triggering timestamp.
A supplied alarm or a local timeout is not proof.
\item For either Force, obtain the captured source context, old and
new versions, retained prefix and boundary from that unique replay
transition. Require $\beta'=\beta_{\rm entry}+1$, $r'=r_b+1$ and
anchor equality with pre-payload $G$, including prefix provenance.
During an unassigned interval, verify the stored context and any
intervening Forces rather than inventing an empty export session.
A later Force supersedes this boundary for future activation, but does
not invalidate its historical cause, original selection or retained anchor.
\item Recompute the first eligible $q_s$ from $q_0$ and all required
historical snapshots. Check target region and immutable next authority.
Set $q^{\rm open}_{\beta',r'}=q_s$ and
$\lambda^{\rm open}_{\beta',r'}$ to the maximum of the anchor's L1
reference height and the initial authority's publication height.
Validate the anchor's original direct or covering commitment evidence
under \cref{sec:coverage}, its complete block and required ancestry.
Do not relabel it with the target pair.
\item Compare the derived fixed fields with $P$ and any existing
fixed record for its target pair. Equal duplicates share that record;
different fixed fields for the same pair are rejected. A packet cannot
replace an already installed handoff. This step records historical facts,
not a claim that the target may still enter.
\end{enumerate}

\paragraph{Permission to enter now.}
\label{sec:can-enter}
$\tagg{CanEnter}(S,P)$ takes the manager's current verified public state
$S$ and a packet whose historical check has passed. It requires that the
target version is current and open, the packet is the still-applicable
fixed boundary, this replica's region is selected, and the local record
has no prior initial entry, close request or Terminated record. Ready
also checks the exact anchor state and dependencies below. A duplicate
already-entered record can be serviced without another Enter.

All initial entry paths use both predicates. Receipt, ancestry and
retained-anchor validation use historical validation only, and global
EXPORT acceptance independently checks its current append permission.
For example, after Force closes version 7, the original boundary and
receipt for retained DATA 108 of $(7,20)$ remain valid history.
$\tagg{CanEnter}$ rejects reopening $(7,20)$, while a valid new Force
packet may inherit that same DATA in $(8,21)$. A cache miss follows the
same historical checks as a cache hit; no prior cached verdict is assumed.

A Force without an eligible successor still has a fixed cause,
anchor, source context and new version. The interpreter retains that
pending boundary. It has no complete selected packet and cannot make
a program Ready until the first eligible historical position fixes
the target region. A later Force can supersede it while preserving the
same stored predecessor context. The synthetic genesis boundary is
verified against the initial checkpoint and selection rule of
\cref{sec:l1}; it does not claim a nonexistent REGULAR, MSIG or AUTO
cause.

\begin{definition}[\defterm{Ready}]
\label{def:ready}
A local Waiting program for $(\beta,r)$ is Ready precisely when:
\begin{enumerate}[label=(\roman*)]
\item historical handoff verification fixes this target pair and selects the
replica's region, or the authenticated genesis rule does so; the separate
current-activation predicate also passes;
\item its exact anchor's original producing pair, height, hash,
complete DATA and commitment evidence have been validated;
\item the native execution state needed to continue after this
anchor has been recovered and verified against that same anchor;
\item the matching finalized authority, fixed quota, separate
reference lower bounds and all validation dependencies are available;
a finalized Espresso header strictly after $q^{\rm open}_{\beta,r}$ is
available with complete verification evidence for the eventual terminal
selection reference;
\item the manager's processed public history has not closed or
superseded the handoff; and
\item the durable record has never performed initial Enter and has
no Terminated tombstone.
\end{enumerate}
\end{definition}

State recovery uses the existing native validation interface. Cached state
for DATA 110 cannot make a program anchored at 112 Ready. A future Force
announcement still permits only Waiting/preparation. The complete regional
message interface below states exactly what is relayed and what its
receiver can conclude.

% Proposed v3.3 message contract. Insert beside sec:ready; no protocol renumbering.
% This fragment does not redefine def:ready or alg:bootstrap.
\paragraph{A complete regional ReadyBundle.}
\label{sec:ready-bundle}
The following is a proposed message interface for the relay used in
\cref{sec:ready}. It makes the transferred state explicit. It is not
a new READY certificate and does not change the public opening of an
export duty. Its timing claim has the precise scope stated below and
requires independent review.

For target session $(\beta,r)$, let $A=A_{\beta,r,0,0}$ and
$\Delta=\Delta_{i_{\beta,r}}$. A \tagg{READYBUNDLE} is an authenticated
regional message with the following typed fields. Every required object
is present as bytes, once, in its normalized dependency set. A URL,
state root, unresolved hash or promise to fetch later is not a substitute.
\begin{table}[htbp]\centering\small
\begin{tabularx}{\linewidth}{@{}>{\raggedright\arraybackslash}p{.25\linewidth}Y@{}}\toprule
Field & Required value\\\midrule
Domain and target & Chain and protocol domain, encoding/native-state
format identifier, target $(\beta,r)$, selected region and sender identity.\\
Fixed handoff & The complete handoff of \cref{sec:witness}, including
its cause, predecessor/source context, selection and fixed target.\\
Exact anchor & $A$'s original producing pair, height, hash, full DATA,
original direct or covering commitment evidence and required ancestry.\\
Executable state $S_A$ & The full, already constructed application-state
object from which execution after $A$ can continue. It includes all state,
code and continuation metadata required by the native execution interface;
if partial execution is supported, its exact continuation position is
included. It contains no imported regional locks, votes or entry timers.\\
State evidence $\Pi_A$ & Native validation evidence connecting this exact
$S_A$ to this exact $A$, with all proof inputs and dependencies. Use the
existing sound application-validation interface; a sender's assertion
or a bare state hash is insufficient.\\
Authority and bounds & Finalized $\tagg{AUTH}_{r,0}$, fixed configuration
and quota, $Q^0_{\beta,r}$, $q^{\rm open}_{\beta,r}$ and
$\lambda^{\rm open}_{\beta,r}$, and all authority evidence needed to
validate the inherited history.\\
Public evidence & An authorized replay checkpoint and complete finalized
L1 control/authority and Espresso header/namespace histories needed to
verify the handoff, references and supplied authorization context. The
checkpoint's trust/authorization dependencies are included. The packet also
contains a finalized header strictly after the opening selection height,
with the evidence required for the terminal reference.\\
Dependency set & Full bytes for the preceding fields' transitive
validation dependencies. A receiver need not fetch an object, replay
application execution to reconstruct $S_A$, or obtain a new authority
before processing this bundle.\\\bottomrule
\end{tabularx}
\end{table}

The native interface determines how $\Pi_A$ verifies state at $A$.
If the block format does not directly commit to a state root, the
bundle supplies the execution evidence and authenticated starting state
required by that interface; this contract does not invent such a
commitment. A merely serialized snapshot that still needs restoration
or conversion is not yet the directly executable state required here.
Recipients independently validate the packet. They do not trust the
sender to authorize a handoff or assert state correctness.

\paragraph{Send and receive events.}
Each correct selected member constructs its complete bundle before its
first Ready/Enter event for this pair. Gathering histories, cross-region
fetching, state reconstruction and bundle formation precede that event;
the existing model assigns them no fixed time bound. At the event it
installs its verified state and executes the existing initial Enter
once, and sends the complete bundle to every member of the fixed
selected-region snapshot. Every correct member does this at its own
first event, so the first correct entrant necessarily sends without
having to know that it is first. Bundle construction or asynchronous
dispatch is not deferred until after the recorded entry time.

The receiver processes the bundle through its serialized Session Manager:
\begin{enumerate}[label=(\roman*),itemsep=2pt]
\item Authenticate the domain, sender and target, and inspect supplied
evidence under the existing evidence-first rule. Validate the historical
handoff, exact anchor, all required authorities, references and native
state evidence. Incomplete or invalid material does not authorize entry.
An existing compatible cache can eliminate duplicate validation, but the
sender does not rely on an unknown receiver cache to omit required bytes.
\item Apply the existing live-activation check against the receiver's
processed public state. A closed or superseded target, a conflicting fixed
handoff, a close request or an exit record prevents initial Enter. This
packet cannot replace a fixed anchor or override newer verified closure.
If the same pair already entered, retain useful evidence without resetting
its state or entry time.
\item For an unentered valid target, validate and adopt the received
$S_A$, check every remaining condition of \cref{def:ready}, then execute
the existing Ready/Enter step of \cref{alg:bootstrap} in this same
message-processing event. The durable entry flag, fixed state attachment
and entry time are recorded atomically before signatures are dispatched.
Relay the complete bundle as above. No separate Ready queue, state-fetch
task or post-Ready installation job lies between this event and Enter.
\end{enumerate}

Here \emph{adopt} has a specific logical meaning: the validated received
state object becomes this new program's executable application state.
The program operates on that object; it does not first reconstruct a
different database from it. This is a representation and message-handler
contract, not a claim that arbitrary database import, disk copying or
native-runtime initialization takes zero time. No receiver can pass Ready
while required installation work remains unfinished.

\paragraph{The exact timing implication.}
Let $e$ be the first correct initial Enter time for this pair. The send
rule places a complete correct-to-correct bundle send at $e$. Suppose the
fixed handoff remains activatable throughout the interval in question
and no closure or other valid entry prohibition intervenes.
If validation, adoption and the atomic Enter transition are the complete
message-processing operation covered by \cref{ass:network}, then
\[
 e_v(\beta,r,0,0)\le e+\Delta
 \qquad\text{for every correct selected member }v.
\]
Indeed, the first correct sender supplies the full executable object and
all validation dependencies. By the existing delivery-and-processing
bound, each correct recipient finishes that handler by $e+\Delta$.
Its guarded final step is Enter, or it has already entered earlier.
There is no remaining state-restoration premise in this argument.

This is the supported abstract-message interpretation, not an additional
service bound for state installation. The quoted network premise covers
delivery and processing of complete validation dependencies. It does not,
by itself, bound an external native import or storage operation merely
because the protocol calls that operation ``installation.'' If such work
remains after bounded message validation, let $u_v\ge0$ be its actual
residual delay. The available decomposition is only
\[
 e_v(\beta,r,0,0)\le e+\Delta+u_v,
\]
when entry eventually completes and remains authorized; the frozen model
gives no uniform bound on $u_v$. Calling that work one event cannot prove
$u_v=0$. In that implementation, the $e+\Delta$ claim and the bootstrap
bounds derived from it remain unsupported.

Consequently the message contract repairs the missing protocol interface
at the abstract level only when the native state is genuinely an executable
validation dependency adopted by the modeled handler. A concrete native
adapter must demonstrate that it meets this contract; sound state proofs
alone do not demonstrate a timing property. The document must withhold
an unconditional native-installation claim until that obligation is met.
No claim is made about an arbitrary history catch-up outside the complete
packet or about when any member first constructs such a packet.

ReadyBundle is regional transport of existing state and evidence. It does
not add a quorum of READY signatures, an Espresso transaction, a publisher
guarantee or a public activation certificate. It neither starts nor resets
an export window. The public clock and the separate live-activation checks
remain those of \cref{alg:force,alg:external,sec:ready}.


\subsection{Activation and bootstrap}
\label{sec:ready}
\begin{protocol}[H]{alg:bootstrap}{Enter and bootstrap the selected session}
Replica $v\in V_{i_{\beta,r}}$ runs these steps for one fixed target
$(\beta,r)$. Let $\Delta=\Delta_{i_{\beta,r}}$.
\begin{enumerate}[label=\textbf{\arabic*)},itemsep=3pt]
\item \textbf{Enter.} In the Ready event, atomically check and set
the durable initial-entry flag. Install $A_{\beta,r,0,0}$,
$\tagg{AUTH}_{r,0}$, $\TB_r$ and both opening reference bounds.
Set the entry lock to $Q^0_{\beta,r}$. Enter $(\beta,r,0,0)$
unfrozen, change Ready to Active, record $e_v(\beta,r,0,0)$ once,
and start the original progress checks. Send STATUS with this lock
and the fixed handoff to the initial leader, and process buffered
proofs for this session. Persist entry before dispatching signatures.
\item \textbf{Choose the anchor.} The initial leader waits
$2\Delta$ from its own entry. If the view remains current and
unfrozen and the retirement contract permits the action, send NEWVIEW
choosing exactly $A_{\beta,r,0,0}$, justified by $Q^0_{\beta,r}$
and the fixed handoff.
\item \textbf{Certify and produce.} Members apply the original
ANCHOR-vote step of \cref{alg:recovery}. On the ANCHOR QC, the leader
starts \cref{alg:steady} at DATA height $h(A_{\beta,r,0,0})+1$
and physical position 1. Local failures use the same recovery rules.
All these steps occur in Active, except when the original Quit,
Completion or explicit retirement transition intervenes.
\end{enumerate}
\end{protocol}

Let $e=\min_v e_v(\beta,r,0,0)$ over correct members. Under the
complete-handler timing contract just stated, if the handoff remains
activatable and bootstrap is not interrupted, all correct members enter
by $e+\Delta$. With a correct initializer and no Quit, NEWVIEW is sent
by $e+3\Delta$, its ANCHOR QC reaches the leader by $e+5\Delta$, and
first DATA reaches all correct members by $e+6\Delta$. These bounds
start at first actual Enter, not at public assignment or candidate creation.
If native installation remains outside bounded message processing, these
bounds have not been established; the extra delay is not silently absorbed
into $\Delta$. No public deadline is reset by either implementation.

\paragraph{The public export clock.}
Public replay assigns an export duty independently of a member's
phase. Its window starts at the containing EXPORT's timestamp for
an already selectable normal successor, at the triggering block's
timestamp for an immediately selected forced successor, or at the
first eligible selection block when the boundary previously had no
assigned duty. A selected program may still be Waiting while that
public window runs. Download completion, Ready, process restart and
duplicate handoffs never reset its deadline. Only an unassigned duty
has no DATA window. Missing evidence pauses replay at its canonical
position; it does not create a later public opening time.

\subsection{Concurrent candidates and races}
\label{sec:runtime-races}
Suppose A's session is $(7,20)$ with terminal DATA 110, B already
produces in $(7,21)$, and $G$ ends at 108. A future MSIG targets
version 7. B can retain its old Active program while preparing
candidate $(8,21)$ from 108. If $G$ advances to 109, that preparation
can change to 109; the old Active program's initial anchor remains
110. If Force then fixes B as successor with anchor 109, all further
preparation is for that fixed anchor. B can enter $(8,21,0,0)$ while
$(7,21)$ closes under \cref{sec:retirement}. Their locks, votes,
entry times and timers are distinct.

If A's terminal 110 enters $G$ before Force, the export context can
advance to B. The candidate then becomes $(8,22)$ at the next eligible
region after B. With cyclic order A, B, C, D, a prospective OFF(C)
that has become effective before the Force selection can redirect
this still-unfixed candidate to D. If Force finally retains B's 112,
the fixed target is D's $(8,22)$ with original anchor producer
$(7,21)$. D must verify state at 112 before starting DATA 113.
Earlier preparation of $(8,21)$ or of state at 109 grants no entry.
OFF cannot rewrite an already fixed regular handoff.

A terminal receipt received before Force can prepare or start a
regular successor in the old version. Once the same node verifies
Force, it cancels an invalid unentered target or requests retirement
of the entered one. If Force arrives first, the late receipt is
retained as regional evidence and cannot resurrect the closed target.
Nodes can briefly have different speculative programs; the same
verified public prefix fixes the same public state and Force boundary.

Every fetch callback retains its original target and anchor identity.
A reply for $(7,21)$ at 110 cannot make $(8,21)$ at 108 Ready.
An old candidate reply may be cached and independently validated for
another use, but it cannot overwrite a fixed anchor. Duplicate receipts,
packets and callbacks cannot run initial Enter twice. A stale MSIG
following an earlier AUTO neither closes the new version nor applies
its attached configuration. Reconciliation discards its obsolete
preparation basis without invalidating a different fixed handoff.


\begin{figure}[p]\centering
\includegraphics[width=.98\linewidth,height=.87\textheight,keepaspectratio]{figures/cross_region_rotation_v3_2/candidate_ready.pdf}
\caption{A changing Force candidate while the old program continues.
The finalized boundary fixes the target and exact anchor; only then can
verified state at that anchor enable Ready and Enter.}
\label{fig:runtime-candidates}
\end{figure}

\subsection{Graceful close and termination}
\label{sec:runtime-close}
Public closure and program exit are separate events. Protocol~5
immediately removes a closed version's permission for new global
appends. Protocol~3 then delivers \texttt{RequestClose} to its local
entered programs, using the verified closure cause. The request is
idempotent and records \texttt{close\_requested}; it is neither Quit
nor terminal Completion. The exact restrictions on new work, treatment
of existing timers and messages, and admissible \texttt{FinishClose}
condition are defined by \cref{sec:retirement}. The drain threshold is derived from the existing regional bound;
resource transfer has no new time bound. A new
session need not wait for the old program to finish closing.

For an unentered program, definitive invalidation cancels its pending
entry and tasks and records Terminated. It has no regional timers or
votes to drain. For Completed, the original Completion action has
already stopped the core. For an Active or Frozen program, only the
retirement contract can authorize \texttt{FinishClose}. An empty local
timer list, a different current public version, or a locally observed
terminal block alone is not that authorization.

In every case resource release requires durable exit and responsibility
records. Required committed DATA, receipts, regional evidence, anchors
and recoverable state must remain accessible under the existing
retention rules. Outstanding complete results and publication retries
are transferred to the Publish Layer, and pending requests in discarded
suffixes remain available for reproposal. Storage and publication
workers acknowledge the responsibilities they accept before the
program releases its last required copy. Later messages can update
those services without reactivating the logical engine. Terminated
does not undo regional commitments, remove DATA from $G$, or authorize
stake withdrawal.

\paragraph{The retirement contract.}
\label{sec:retirement}
Fix a session $(\beta,r)$ and write $\Delta=\Delta_{i_{\beta,r}}$,
using its immutable voting snapshot. Retirement adds a local
\texttt{close\_requested} flag, closure evidence, and a first-close time
$c_v$ to the program record. These are not another phase. While draining,
an Active or Frozen program keeps that phase; only a genuine Quit or
terminal Completion changes it. The following rules extend the running
core explicitly. They do not make Force a quit proof or a terminal receipt.

\paragraph{Complete-receipt dissemination.}
\label{sec:retirement-receipt-relay}
The following obligation applies during normal running as well as
retirement. The first time a correct regional member or its local
manager/collector holds a complete valid receipt of its own session,
it multicasts the receipt and all validation dependencies to the
session's fixed membership. In particular, a terminal receipt is
multicast no later than its local Completion event. Deduplication by
session and DATA hash suppresses repeated encodings or signer sets,
not this first complete relay. The relay and any pending delivery are
preserved across Completed and Terminated by the persistent owner.
Recipients process an own-session terminal receipt under Completion
even while Frozen or closing, independently of its generation and
before a simultaneous timer. A program that has not entered caches
it; a Terminated program archives it without restarting. This rule
forwards existing signatures and introduces neither a new vote nor a
certificate. The intra-region complete-message bound applies to the
actual relay; it is not a bound on first forming a receipt.

\paragraph{\texttt{RequestClose}.}
The Session Manager applies this operation once per program as follows.
\begin{enumerate}[label=(\roman*),itemsep=2pt]
\item Verify complete evidence that the canonical prelude of an
authenticated finalized Espresso block closed this program's version
$\beta$. The evidence contains the triggering header, the applicable
complete control and namespace history from a validated checkpoint,
and the resulting closure and retained-prefix context. Recompute the
existing MSIG/AUTO checks; a pending MSIG, an OFF, an unfinalized header,
or a bare timestamp is insufficient. Closure does not require a selected
or Ready successor.
\item At the first such verification, atomically set
\texttt{close\_requested}, record the current time $c_v$ and the evidence,
and relay its complete validation bundle to the fixed regional membership.
This is transport of existing source proofs, not a regional vote or new
certificate. Persist a closed-version record that also prevents delayed
witnesses from creating or entering another old-$\beta$ program.
Duplicates retain the original $c_v$.
\item Close every fresh-action gate for this instance: no DATA/DRAIN
proposal, NEWVIEW, DATA/DRAIN/ANCHOR vote, new commit timer, progress
check producing BLAME, term entry, or generation entry is allowed.
The gate covers initial Enter, delayed authority/recovery callbacks,
terminal-QC rotation and buffered-message processing. No new STATUS
or other entry action can bypass it. Already dispatched messages are
not withdrawn. New-session preparation and entry use their own pair
and need not wait for this instance to retire.
\item Keep existing commit timers and their original deadlines.
Continue complete validation, comparison with retained signed material,
and evidence-first processing. A genuine current-generation quit proof
still runs Quit, broadcasts the proof, cancels its timers and preserves
commitments; it cannot reopen the entry gate. A surviving timer still
commits its block and ancestors. For each DATA thereby committed, its
Protocol~1 FINAL is delivered both to the fixed regional membership and
to the publication collector, preserving the regional dissemination rule
of \cref{def:receipt}. A partial FINAL does not invoke the complete-result
guarantee of \cref{ass:publish}.
\item Continue collecting old votes, QCs, FINAL messages and receipts.
Apply the complete-receipt multicast rule above, and offer every
complete result to the Publish Layer. A valid terminal
receipt still executes genuine Completion, possibly changing the phase
to Completed and canceling timers; a pending close retains its drain
deadline. Other valid receipts retain their original processing rules.
Neither receipt arrival nor a later ON reopens the closed global branch.
Keep data service and publication duties until their durable transfer.
\end{enumerate}
Messages, including valid quit evidence, are processed before a
simultaneous timer expiry or retirement event. No receipt, lock, vote,
committed block or provenance field is erased by RequestClose.

\paragraph{\texttt{FinishClose}.}
For a program that was Active or Frozen when RequestClose ran, set
\[
 d_v=c_v+4\Delta.
\]
FinishClose may mark it Terminated only at a time $t\ge d_v$ when:
\begin{enumerate}[label=(\roman*),itemsep=2pt]
\item every existing commit timer has executed or was canceled by
genuine Quit or Completion, and no queued action can pass a fresh-action
gate;
\item persistent storage owns the exit and closed-version records,
committed log, required signed evidence and validation dependencies,
receipts, anchor data and recoverable state, without releasing the last
required copy;
\item a persistent owner has accepted outstanding complete-receipt
multicasts, complete-result offers, publication retries and data-service
duties, and each locally complete result has been offered to the
Publish Layer; and
\item the dispatcher routes later messages for $(\beta,r)$ to those
owners. They may collect or relay existing signatures, validate late
receipts, serve data and continue publication, but cannot restart the
engine, execute its BFT timers or issue fresh BFT signatures for it.
\end{enumerate}
The ownership transfer may be local to the same correct replica; it
requires no additional agreement or certificate. Only after these
guards hold may the production/recovery worker and its timer machinery
be released. A late terminal receipt records archival evidence and does
not change a Terminated program back to Completed or Active.

A definitively invalid Waiting or Ready program has never entered
and has no regional timers to drain. A program already Completed when
retirement is requested has performed the original timer cancellation.
These programs may terminate after the corresponding durable exit and
responsibility transfer without the $4\Delta$ drain. A temporary change
of an unfixed candidate's eligibility alone need not terminate it.
Termination never deletes a commitment, removes DATA from $G$, or
authorizes stake withdrawal.

\begin{proposition}[Closure propagation and honest-message drain]
\label{prop:retirement-drain}
Under \cref{ass:network}, let $c=\min_{v\ {\rm correct}}c_v$ be the first
correct RequestClose time in this region, with the complete bundle
specified above. Then every correct member has closed by $c+\Delta$.
Every surviving honest commit timer expires by $c+3\Delta$, and the
regional FINAL messages generated by those timer events reach every
correct member by $c+4\Delta$. Consequently $d_v=c_v+4\Delta$ is a
conservative local drain threshold, not a deadline for finishing all
retirement duties.
\end{proposition}
\begin{proof}
The first correct relay is a complete correct-to-correct message.
Its delivery and validation are covered by the existing regional bound,
so $c\le c_v\le c+\Delta$. A correct DATA/DRAIN vote at time
$s=\tau_v(P)$ starts its only timer, with expiry $s+2\Delta$.
The close gate gives $s\le c_v\le c+\Delta$; no callback starts a
replacement timer. Thus its expiry is at most $c+3\Delta$.
Quit and Completion can cancel a timer but cannot postpone it.
The FINAL messages actually generated at surviving expiries are
delivered regionally within another $\Delta$. Finally $c_v\ge c$
implies $d_v\ge c+4\Delta$. Message-first processing covers equality.
\end{proof}

The bound starts at verified closure and its complete regional relay,
not at a control signature, the triggering Espresso timestamp, or the
start of unbounded historical catch-up. It does not add a new source
visibility assumption. Three local $\Delta$ intervals cover the last
honest timer event but not its FINAL delivery: a member can vote just
before $c+\Delta$, commit just before $c+3\Delta$, and deliver FINAL
only just before $c+4\Delta$.

Four $\Delta$ intervals do not bound subsequent receipt relays,
Byzantine message release, a storage copy, $T_{\rm pub}$, $T_E$, or
successful EXPORT acceptance. Their outstanding duties pass to the
persistent owners; no bounded storage size or transfer latency is
claimed. Elapsed grace alone therefore never authorizes FinishClose.

\begin{proposition}[Conditional regional safety during retirement]
\label{prop:retirement-safety}
Assume the stated regional model, the baseline regional validation
and DATA/DRAIN soundness premises, and the complete-message Ready/Enter
installation contract of \cref{sec:ready}. With the complete-receipt
multicast and retirement rules above, correct regional committed logs
remain prefix-comparable and irrevocable. The external prefix argument
of \cref{sec:safety-argument} is preserved. This does not assert a
protecting QC for every late surviving timer, old-session progress,
or complete-receipt formation.
\end{proposition}
\begin{proof}
We verify the prerequisites of the timer/lock argument explicitly;
baseline safety is not used to infer an unproved universal-QC property.
Fix a correct direct timer for DATA or DRAIN block $B$ in generation
$g$, voted at $s$, that actually expires at $s+2\Delta$ without
cancellation. All comparisons below concern this same session.
Let $e_0$ be its first correct initial-entry time. The specified Ready
contract enters every nonclosed correct member by $e_0+\Delta$.
An own-session terminal receipt requires a correct FINAL after a timer,
whose earliest possible expiry is $e_0+2\Delta$. Hence such a receipt
cannot evade Completion by arriving only at an unentered cache of a
correct member that will run the chosen timer.

\emph{Quit and Completion cannot suppress its early vote window.}
If any correct member first quits $g$ at $a\le s+\Delta$, its complete
quit proof reaches the committer by $a+\Delta\le s+2\Delta$ and
cancels the timer. Thus any first correct Quit satisfies
$a>s+\Delta$. If any correct member Completes at $b\le s+\Delta$,
its complete terminal receipt is multicast by $b$ and reaches the
committer by $b+\Delta\le s+2\Delta$. The committer then executes
Completion: it commits the terminal DATA and ancestors and cancels
this timer, including a DRAIN timer. Again this contradicts the chosen
surviving direct timer. No correct Completion can therefore occur at
or before $s+\Delta$. This is an operational cancellation argument;
it does not assume that receipt-triggered commitments are already
known to be prefix-comparable.

\emph{The proposal is processable within its regional window.}
For $g=0$, the self-contained Ready relay and installation contract
in \cref{sec:ready} supplies entry by $s+\Delta$; this assertion is
conditional on that precise contract, not a bound on state restoration
before any member is Ready. For $g>0$, let $a^-$ be the first correct
Quit of $g-1$, and let $e$ be the first correct entry into $g$.
The original recovery rule gives $e\ge a^-+\Delta$ and $s\ge e$.
The complete quit relay reaches every correct member by $a^-+\Delta$;
its own mandatory wait has therefore ended by
$a^-+2\Delta\le e+\Delta\le s+\Delta$. The forwarded complete
proposal supplies the finalized authority and remaining admission
history. Hence a receiver with no Quit, Completion or close preventing
entry can process it in $g$ by $s+\Delta$. Applying the same
complete-proof argument inductively to earlier generations supplies
any predecessor entry needed to process the quit proof. This uses
the existing recovery waits and complete validation bundles, not a
new simultaneous authority-visibility time.

\emph{Early and late windows.}
Let $c$ be the first correct RequestClose time. If $s+\Delta<c$,
no close, Quit or Completion suppresses any correct member's vote
while $B$ propagates. Earlier-term processing and prior indirect
commitment do not suppress that vote either. An incompatible retained
history would expose a genuine quit proof in this window and contradict
timer survival. Every correct member therefore has cast, or now casts,
the matching DATA/DRAIN vote by $s+\Delta$. Their total weight is at least
$W-F>F$, so the votes form $B$'s protecting QC by $s+2\Delta$ at
all members that can still take a later recovery entry. A member that
Completed after the vote window cannot take such an entry; it need
not be assigned an unused recovery lock.

Any correct entry into $g+1$ must obey
\[
 e_v(\beta,r,g+1,0)\ge a_v+\Delta
 \ge a+\Delta>s+2\Delta.
\]
In the early case, every such entrant therefore knows the protecting
QC before its lock snapshot. The highest-QC rule gives the baseline
recovery-lock invariant. Closure preserves collecting those already
sent votes and QCs, and may only remove a later entry. In the late
case $s+\Delta\ge c$, the timer expires at or after $c+\Delta$.
Every correct entry gate has closed by then; the strict displayed
inequality prohibits every later-generation entry and thus its
correct signatures. A later-generation QC needs a correct signer,
so no recovery can exploit a missing protecting QC in this case.

\emph{Direct timers and receipt-based commits.}
In one generation, complete signed histories either extend one
another or expose the original admissible equivocation proof.
A correct voter forwards its complete history. A correct signer of
an incompatible history either processes $B$ first and cannot vote,
or forwards its history before processing $B$, no later than
$s+\Delta$. That history reaches the committer by expiry and exposes
the proof. Genuine Quit, or Completion before expiry, would cancel
the timer. Thus no conflicting same-generation correct direct timer
can survive; the same argument excludes a conflicting QC, including
one assembled later from old votes.

Across generations, a surviving lower-generation direct timer cannot
be in the late case if a correct higher-generation vote exists.
Its early-case QC is known before every later live lock snapshot.
The baseline ranked-QC descent now has its actual prerequisites:
correct ANCHOR voters meet those locks; DATA/DRAIN QCs inherit a
lower-ranked parent justification. It excludes a later QC that does
not extend the committed block, and the parent of a later direct
proposal consequently extends it. These are precisely the
protecting-QC, recovery-lock and unique-extension steps of Part~I,
with the explicit vote-window checks above and the baseline
DATA/DRAIN validation rules. Direct commitments are comparable, and
ancestor commitments inherit that property.

Finally a valid receipt contains a correct FINAL signer. That signer
sent FINAL only after a timer committed its named DATA, directly or
as an ancestor; committing from another receipt emits no FINAL.
Consequently every receipt-based commit, including Completion, traces
to a causally earlier direct timer commitment of that DATA or a
descendant. It adds an already protected prefix. This provenance
argument does not take the safety of receipt-based commits as its own
premise. A timer canceled by Completion contributes no additional
direct commitment, particularly no commitment of its unexpired DRAIN.
No retirement step removes a prior commitment.

The public interpreter still preserves all of $G$ at its canonical
Force boundary and rejects further old-version DATA. Retirement and
receipt relay neither select a retained anchor nor authorize an
append, so those external transitions are unchanged.
\end{proof}

The universal protecting-QC lemma from Part~I cannot simply be copied:
its proof requires matching votes from all correct members, whereas a
closed member now retains a proposal without voting. For example, with
three unit-weight members and $F=1$, let correct $H_0$ close at $c=0$
and correct $H_1$ receive closure at $\Delta$. A Byzantine leader gives
a valid proposal to $H_1$ at $0<s<\Delta$. $H_1$ votes, forwards and
later commits; $H_0$ retains it without voting. If the Byzantine member
withholds its vote and FINAL, only $H_1$'s FINAL exists. Its weight is
not greater than $F$: neither a proposal QC nor a complete receipt need
ever form. Retirement preserves this local commitment and evidence
without calling it Completion. The Publish Layer's complete-result
premise does not repair the missing signature.


\section{The Force handler}
\label{sec:force}
Protocol~5 closes the external branch. Protocol~3 consumes the verified
closure and separately retires old programs under \cref{sec:retirement}. The interpreter derives an automatic Force itself
or applies a previously announced multisignature Force.
Both retain the same object: the entire global prefix immediately
before the triggering Espresso block's payload.

\subsection{Export duties and block-missing windows}
\label{sec:windows}
An export duty belongs to the session responsible for the next DATA
height after $G$. Its window parameters come from its fixed initial
authority, not from a later parameter update.
Let $R$ be its remaining DATA quota when a window opens. All targets
below are positive because completion removes the duty before $R=0$
can open another window.

At public opening timestamp $\theta_0$, create the startup window
\[
 [\theta_0,\theta_0+W_{\rm start}),\qquad
 m=\min\{m_{\rm start},R\}.
\]
After it succeeds, use consecutive fixed windows of length
$W_{\rm run}$, each with target $m=\min\{q_{\min},R\}$ measured from
that window's opening. The window stores its start, deadline, target,
accepted count and phase. A new block increments the count only after
its EXPORT is accepted into $G$. Each DATA is counted once in its own
session. Empty DATA counts under the same quota rule; DRAIN, heartbeats,
proposal QCs, duplicate records and future-session data behind a gap
do not count.

At the first Espresso block crossing the deadline, compare the count
with the target. Success opens the next fixed window and resets its
count. Extra output is not banked as unlimited credit for later windows.
If several deadlines have passed, settle the successful windows in
order and stop at the first failure. It produces one automatic Force.
A terminal DATA accepted before its deadline cancels that session's
remaining checks, even if the final window has not expired.

The intervals are half-open. DATA with a timestamp equal to the
deadline is too late for that window. The check runs before the first
block's payload at that timestamp, not after all blocks that may
eventually share the timestamp.

\paragraph{What a failure establishes.}
Failure means that the next required output was not delivered to the
global prefix at the specified rate. It does not identify a faulty
proposer, submitting participant, region or network. Failure to form a
complete receipt, restore state, deliver a result, or obtain Espresso
service can cause the same observation. A missed service bound violates
that liveness premise; the Force rule does not identify the cause. It is a recovery condition,
not automatically slashable equivocation or a proof of censorship.

\subsection{Scheduled control and two kinds of Force}
\label{sec:force-stages}
An MSIG Force binds the chain, target branch version $\beta$, activation
timestamp $\theta_a$, control sequence and authorized configuration
changes. Its L1 sequence identifies the instruction. It contains no
retained DATA height. The committee submits it under
\cref{ass:committee,ass:monitor} and serves its finalized evidence.
Signing alone does not start the submission bound.

L1 finalization makes the instruction known, but does not yet close
the branch. Its effect begins only at $J_{q(\theta_a)}$.
Until then, normal production, normal handoffs and valid exports can
continue. A node does not veto future global output merely because a
Force has been announced.

At activation, execute the instruction only if its target version is
still current. If an earlier Force changed that version, mark the
instruction processed without closing the new branch or applying its
attached configuration. A normal regional handoff does not change
$\beta$, so it does not cancel the instruction. A separately scheduled
STATUS remains independent. These target checks and the priority of
MSIG over AUTO are unchanged from version 3.

An automatic Force binds the chain, old branch version, export-session pair,
failed-window identity and triggering Espresso header. Its cause
deadline is the failed window's deadline; its realized time is the
timestamp of that triggering block. No transaction is needed to create
this derived event. An optional published audit record does not grant
it a second effect.

\begin{protocol}{alg:force}{Process the prelude and close a branch}
Before processing the payload of each authenticated finalized Espresso
block $J_q$, let $\theta=\theta_q$ and capture the entry branch version.
\begin{enumerate}[label=\textbf{\arabic*)},itemsep=3pt]
\item \textbf{Apply status.} From the complete control prefix, apply
all due independent STATUS and CLEAR records not yet applied, in
$(\theta_a,c)$ order. Mark them applied. Later blocks with the same
timestamp do not repeat them.

\item \textbf{Apply a multisignature Force.} Examine due Force
records in that same order. Apply only the first record whose target
matches the entry branch version. Capture the assigned export duty's context,
or the stored boundary context while no export duty is assigned. Apply that record's
attached configuration and close the branch version. Consume every other due
Force as stale, without applying its attached configuration. If a
matching Force closes the branch version, cancel its old export window and
skip step 3.

\item \textbf{Check delivery.} If an export session exists, settle
its deadlines at or below $\theta$ as in \cref{sec:windows}.
On the first failed window, capture that export duty's source region,
index and branch version,
set its cooldown to $\theta+C_{\rm auto}$, and close the branch version.
An unavailable namespace or control proof cannot establish a failure.

\item \textbf{Preserve and reopen.} For either closure, retain all
of $G$ and its provenance. Archive the old branch version's unconfirmed suffix
and normal successors. Record the cause, triggering header, retained
anchor, captured source context and selection lower bound $q$.
Increment $\beta$ by one. Do not accept new-version DATA in $J_q$ itself.

\item \textbf{Select or wait.} Record the pre-payload eligibility
snapshot $\chi_E(q)$. Resolve the new boundary, or an older unassigned
boundary, by its first eligible historical selection position.
When a duty opens, start a fresh startup budget at its canonical
opening timestamp. If no candidate exists, keep the anchor and run
no DATA window. Then return to payload processing.
\end{enumerate}
\end{protocol}

The cooldown is a replay-derived eligibility restriction. Governance ON
alone does not clear it; an explicit CLEAR can. Governance OFF continues
to exclude a region even after its cooldown expires.
Window and quota parameters remain fixed for the duty. A control that
needs to end it early uses an MSIG Force, not a retroactive deadline
change.

If Espresso time jumps across several scheduled thresholds, all those
events are first observable at the same block boundary. Step 1 processes
the independent status changes before either Force, and the matching
MSIG Force takes precedence over the entry session's automatic check.
No DATA is inserted between these prelude events. New startup windows
begin at the realized selection timestamp, not at an old missed
deadline, so a single block does not time out every new successor.

\subsection{Processing rules and control completeness}
\label{sec:control-completeness}
Control ingestion verifies a full L1 prefix and retains a pending
activation queue. A sequence number alone does not mean all ingested
records have activated. The interpreter processes each control once,
in its recorded target context. A closed branch version is never reopened by
an ON update, a repeated Force or a late regional receipt.

Every correct observer refreshes its complete finalized control view
before confirming an Espresso interval. The advance-announcement
assumptions ensure that all controls applicable to that interval are
already available. Missing evidence pauses replay. A new observer
reconstructs the same history; it does not reject an old instruction
merely because it learned it late.

A certifier's pure core must bind its supplied L1 checkpoint, control
sequence and history commitment. Its L1 acceptance interface must
reject an attestation that omits applicable controls. A conservative
implementation requires the advertised complete control version to
match the version required by the accepting contract; a stale version
is refreshed before retry. More selective checks require an equivalent
completeness proof. Proof of one event's inclusion is not proof that
no intervening control exists.

Discovery of a previously omitted control that should have affected
already confirmed history violates the completeness or announcement
premise. It is not a new permission to roll back $G$, apply the Force
at an arbitrary later point, or let nodes choose by local arrival time.
This protocol's conditional guarantees do not specify recovery from
that premise violation.

\section{The external observer}
\label{sec:observers}
An observer keeps the provisional production branch, persistent global
prefix, source-verification cursors and interpreter state. It does not
store per-block relay clocks or a committee-protected prefix.
The same deterministic logic can run inside the certifier.

\begin{protocol}{alg:external}{Verify Espresso output and advance the prefix}
\begin{enumerate}[label=\textbf{\arabic*)},itemsep=3pt]
\item \textbf{Authenticate inputs.} Obtain the next finalized
Espresso block and complete namespace evidence, linked to the trusted
quorum and prior prefix. Obtain the complete applicable finalized L1
control context and required historical authorities. If any required
history is unavailable, wait without treating it as empty.

\item \textbf{Run the prelude.} Execute \cref{alg:force} before
reading this block's payload. Persist the eligibility snapshot for this
Espresso position. Notify regional nodes of any derived forced boundary.

\item \textbf{Accept exports in order.} Decode records in canonical
namespace order. For each EXPORT, run \cref{sec:exports} on a temporary
state, including each normal boundary and session count. Commit the
whole update only if it passes. Invalid candidates do not change $G$
or restart a clock. Supporting evidence can be cached after validation.

\item \textbf{Persist and expose.} Persist the resulting state and
source cursors. Newly accepted DATA is global now, without another
timer or L1 transaction. Expose provisional output separately.
On closure, archive only the unconfirmed suffix and keep all of $G$.
\end{enumerate}
\end{protocol}

The expected base of an EXPORT is the DATA frontier and its binding
prefix digest, with provenance. It is not the hash of the entire
interpreter state: an empty Espresso block can change time, controls
and windows without changing this base. The certifier's state-transition
attestation, in contrast, binds the complete previous and next replay
checkpoints as well as the published range.

\subsection{Requests and retained data}
\label{sec:requests}
A request is \defterm{admissible} when it passes the application's admission
checks. Each node queues it in order of first receipt. Repeated delivery
does not change its queue position. New arrivals do not move
ahead of queued requests, and different nodes need not have the same
queue order. A correct leader skips requests
already in its proposed ancestry, including the inherited regional
prefix, and fills the batch with the oldest remaining requests that fit.
If no request can be included, it may propose an empty DATA block.

A request is finally removed from the pending queue only when included
in the global prefix. A request in a discarded regional suffix remains
pending and can be proposed on the new branch. Regional commitment or
regular handoff alone does not remove it. Keep the block data and
receipts needed for recovery and checkpoint verification. Before finalized Espresso
inclusion, keep complete results available to the Publish Layer and
direct readers. After finalization, other regions may retrieve them from
authenticated Espresso history. Preserve required archival or checkpoint
coverage; global promotion does not authorize deleting the only copy.



\paragraph{Faithful publication.}
The proposer or batcher publishes the next complete interval after
$C_{\rm pub}$, even if a Force has changed the current region.
The interval must belong to $G$, preserve order and multiplicity, and
match the DA payload and native execution evidence. An unavailable
accepted payload causes waiting, not an invented replacement.
The exact native adapter is specified separately: a Force cannot move
back from a partial-EVM anchor simply because the next signer cannot
continue it. It must obtain a legal continuation or wait.

\section{Examples and guarantees}
\label{sec:examples}
\subsection{Regular rotation and a late DRAIN}
\label{sec:regular-example}
Let $h(A_{\beta,r,0,0})=100$, $\TB_r=10$ and $K=4$. The normal DATA terms
produce 101--104, 105--108 and 109--110. A terminal QC below 110 moves
to the next term. A receipt for DATA 110 opens the next region, which
starts DATA 111 after Ready and bootstrap. No global wait is inserted.

\begin{figure}[htbp]\centering
\input{figures/sync_hotstuff_rr_rotation_v3_1/regular-handoff.tex}
\caption{Normal cross-region rotation retains the DATA, FINAL and
ANCHOR sequence. The historical Espresso selection snapshot fixes the
successor. Regional production and Espresso publication proceed in
parallel. Neither the terminal receipt nor the new region waits for
global confirmation.}
\label{fig:regular-cross-region}
\end{figure}
\FloatBarrier

Now suppose 110 has a QC but its timer is cancelled by Quit.
Recovery opens at height 110 and proposes a DRAIN, which supplies a
new timer without using DATA quota. When it commits the terminal DATA,
replicas send FINAL for 110, not for DRAIN. The normal successor still
inherits 110. Its old L1 and Espresso references do not establish an
earlier commitment or export time.

Suppose $G$ is only at 100 while A has committed 110 and B has committed
111. The Publish Layer can export 101--111 with the required A and B receipts
and normal witness. Once that complete EXPORT is accepted from finalized
Espresso history, the entire interval becomes global. A lone record for
111 cannot bypass missing ancestors.

\subsection{A multisignature Force and the retained anchor}
\label{sec:force-example}
Let the region order be A, B, C. A's session $(7,20)$ produces
DATA 101--110. Its terminal receipt lets B start session $(7,21)$
before Espresso has confirmed all of A's DATA. Suppose B has produced
through 115, while $G$ ends at 108. A is still the export region:
its session owes the next DATA, 109.

The committee submits an MSIG instruction targeting version 7 with
activation time 120 and an attached OFF(A). It submits early enough
for the full L1 visibility bound. Before the first finalized Espresso
block with timestamp at least 120, valid old-version exports may still
advance $G$.

Suppose $G$ still ends at 108 at that boundary. Nodes and certifiers
check that the current version is 7, apply OFF(A), and close version 7
before reading the block's payload. They preserve 108 and reject any
later export of the old A:109--110 and old B:111--115. The old regional
receipts remain evidence of regional commitment, not permission to
append to $G$.

If B is eligible, it opens a new session $(8,21)$ from 108. Its entry
view is $(8,21,0,0)$. This is a different engine from $(7,21)$.
After Ready and the normal ANCHOR bootstrap, it proposes DATA 109.
Its new-version DATA can be accepted only in a later Espresso block.
There is no second Espresso-to-L1 confirmation round before bootstrap.

Two variations explain the target check. If AUTO already changed
version 7 to 8 before timestamp 120, the MSIG instruction is stale:
it neither closes version 8 nor applies its attached OFF(A). If A
instead completed its export and B became the export region through
normal handoff, the version is still 7. The instruction still applies,
preserves the newer $G$, and selects after B. To make OFF(A) apply even
when the Force becomes stale, submit it as an independent STATUS.

With a DATA quota of ten, the two B sessions can have different
endpoints: $(7,21)$ starts after 110, while $(8,21)$ starts after 108.
They read the same $\TB_{21}=10$ but have endpoints 120 and 118.
The pair identifies the engine; the numeric index selects its authority.

\begin{figure}[htbp]\centering
\input{figures/sync_hotstuff_rr_rotation_v3_1/forced-handoff.tex}
\caption{Multisignature Force: L1 supplies a future-dated instruction;
the first finalized Espresso timestamp crossing fixes its boundary.
The complete pre-boundary global prefix survives. State transfer and
the original regional bootstrap follow; old regional engines are not
silently stopped.}
\label{fig:forced-cross-region}
\end{figure}

\FloatBarrier
\subsection{Four concrete view updates after Force}
\label{sec:force-view-examples}
The view is $(\beta,r,g,t)$: branch version, regional turn,
recovery number and leader term. An effective Force increments $\beta$;
the new session starts with $g=t=0$. Its turn number is $r_c+1$, where
$r_c$ is the session that owes the next DATA after $G$. It is not
necessarily the latest production turn.

The following examples use the region order A, B, C. Each effective
Force has an eligible successor that completes Ready and the normal
bootstrap. MSIG instructions satisfy the advance-announcement condition.
The first three examples apply to either an effective MSIG Force or
an AUTO Force; the fourth explains why a stale MSIG does nothing.

\Needspace{10\baselineskip}
\paragraph{Example 1: A produces and also owes the next DATA.}
A is in view $(7,20,2,3)$. Its session produces DATA 101--110,
but $G$ ends at 108. A still owes DATA 109, so $r_c=20$.
A Force selects B, whose new entry view is
\[
 (7,20,2,3)\xrightarrow{\text{Force}}(8,21,0,0).
\]
Version 7 becomes 8. Turn 20 is followed by turn 21. Recovery and
leader numbers restart at zero because B opens a new session.
B starts DATA 109 after bootstrap. DATA through 108 is retained.

\paragraph{Example 2: normal handoff and global confirmation have reached B.}
A first hands off normally:
\[
 (7,20,0,2)\xrightarrow{\text{regular}}(7,21,0,0).
\]
B then changes leader and reaches $(7,21,0,1)$. Suppose $G$ now ends
at B's DATA 113. B owes DATA 114, so $r_c=21$. A Force selecting C gives
\[
 (7,21,0,1)\xrightarrow{\text{Force}}(8,22,0,0).
\]
C starts DATA 114. A scheduled MSIG targeting version 7 still matches:
the earlier regular handoff changed $r$, not $\beta$.

\paragraph{Example 3: C produces, but global confirmation still waits for A.}
A's session $(7,20)$ has committed DATA 101--110. B's session
$(7,21)$ has committed 111--120. C is producing later DATA in view
$(7,22,1,2)$. However, $G$ still ends at 108. The next required DATA
is A's 109, so the export turn is still $r_c=20$.

A Force selects after A. If B is eligible, its new entry view is
\[
 (8,21,0,0),\qquad\text{not }(8,23,0,0).
\]
The calculation uses the unfinished export turn 20, not C's latest
production turn 22. Old A:109--110, old B:111--120 and C's later output
cannot append after closure. New B starts DATA 109 from the retained 108.
Although the numeric turn falls from the latest production index 22 to
21, no globally confirmed DATA is removed. Sessions $(7,21)$ and
$(8,21)$ have separate engines.

\paragraph{Example 4: AUTO has already made a scheduled MSIG stale.}
An MSIG instruction targets version 7 at Espresso timestamp 120.
At timestamp 115, AUTO instead causes
\[
 (7,20,2,3)\xrightarrow{\text{AUTO}}(8,21,0,0).
\]
By the first crossing of 120, suppose the new session has reached
$(8,21,1,2)$. The MSIG still targets 7, but the current version is 8.
It is marked processed without executing. Assuming no other transition
is due, the view stays
\[
 (8,21,1,2)\xrightarrow{\text{stale MSIG}}(8,21,1,2).
\]
None of the four fields changes. Its attached configuration also has
no effect; an independent STATUS follows its own activation rule.

\paragraph{If no successor is available.}
An effective Force still increments $\beta$ and retains $G$, but the
public boundary remains unassigned. It has not yet opened a new regional view.
For example, after closing export session $(7,20)$ it waits in version
8; a selected successor can later enter $(8,21,0,0)$ at Ready.
If another effective Force acts before any successor opens, the version
becomes 9 but the stored predecessor index remains 20. The eventual
entry view is then $(9,21,0,0)$, not $(9,22,0,0)$.

\subsection{A failed window and equal timestamps}
\label{sec:window-example}
For illustration, let the startup window be $[0,4)$ with target one,
and later windows have length four and target two, subject to remaining
quota. One DATA at time 1 satisfies startup. If only one further DATA
is accepted at time 5, the running window $[4,8)$ fails.
A record in the first block timestamped 8 is too late to repair it.
The automatic Force preserves the two accepted DATA and closes their
branch version before reading that block's payload.

\begin{figure}[htbp]\centering
\input{figures/sync_hotstuff_rr_rotation_v3_1/observer-finality.tex}
\caption{External confirmation and a block-missing Force. Only accepted
DATA counts. The running-window example assumes at least two DATA
remain due; in general the target is capped by remaining quota.
Deadline-equal output cannot revive the expired branch version.}
\label{fig:observer-finality}
\end{figure}
\FloatBarrier

If Espresso jumps from timestamp 5 to 30, the same first failed window
is detected at the latter block. A new successor starts its window at
30, not at 8. If another Espresso block also has timestamp 30, it does
not repeat the old control events or Force. A due status update at the
same crossing is applied first, so the successor uses the prescribed
eligibility snapshot.

\subsection{Conditional safety and lifecycle invariants}
\label{sec:safety-argument}
The statements below separate global interpretation from regional
execution. They do not assert liveness from safety. The model and
verification interfaces remain those of Section~2.

\begin{proposition}[Common interpretation and preservation of global DATA]
\label{prop:global-prefix}
Let two correct observers start from the same authenticated interpreter
checkpoint and apply Protocols~5 and~6 under the stated history,
announcement and verification premises. Their global DATA sequences are
prefix-comparable. Replaying the same complete Espresso prefix yields
the same interpreter state. Every accepted DATA block retains its content,
position and producing identity forever. Once a version is closed, no
later EXPORT adds new DATA from that version to $G$.
\end{proposition}
\begin{proof}
Induct over the common finalized Espresso sequence and, within a block,
its prelude and ordered EXPORT records. At the initial checkpoint the
states agree. Complete applicable L1 controls, immutable authorities,
Espresso headers and namespace contents authenticate the same next
input. Advance announcement ensures the relevant control prefix is
available before that crossing; a verifier lacking evidence waits.
The status order, first matching MSIG, alternative AUTO check, fixed
windows and first eligible historical selection are deterministic.
They therefore produce the same prelude state.

An EXPORT is tested against that state and the exact current DATA base.
Regional and native validation use the stated sound interfaces and the
same data. Acceptance atomically appends exactly the next continuous
range; rejection appends nothing. This gives identical states after
each record. Different verified prefix lengths are ordered by Espresso's
prefix property, so one observer is at an earlier step of the same
state sequence.

Inspect preservation at each step. STATUS, CLEAR and window settlement
remove no DATA. Force retains all of $G$ before payload and changes only
the authority for future appends. Selection of an unassigned boundary
also preserves $G$. EXPORT only extends it. Retained DATA keeps its
original producing pair and hash. Closed versions are never reopened,
and every later EXPORT checks the current version for new DATA. Thus
neither an old regular successor nor a late receipt can append a closed
suffix. Local Handler phases and retirement do not mutate $G$, so they
add no exception to this induction.
\end{proof}

\begin{proposition}[Fixed entry and isolation of local programs]
\label{prop:program-isolation}
Under the regional-validity and source-verification premises, each
opened pair $(\beta,r)$ has one fixed handoff. At a correct replica,
initial Enter occurs at most once for that pair, uses exactly that
handoff's anchor and state, and never uses an obsolete task callback
as authority. Actions for one pair do not alter another pair's locks,
votes, timers or entry time. Candidate preparation creates none of
these BFT actions.
\end{proposition}
\begin{proof}
Proceed from the initial checkpoint through successive handoffs.
For a regular handoff, the predecessor's fixed boundary and quota fix
its terminal height; regional validity fixes the terminal DATA and its
signed selection reference. The least eligible position fixes the
region and designated authority. For Force, common replay fixes the
old and new versions, source context and exact retained $G$; historical
selection then fixes the target. Repeated Forces without assignment
change $\beta$ but keep the stored predecessor index. They cannot
reuse an already opened pair. Handoff verification rejects conflicting
fixed fields instead of replacing an existing record.

A preparation may have several candidate anchors over time. None is
eligible for Enter until verification installs the actual handoff.
Material callbacks retain the requested pair and anchor identity;
the gate checks the exact fixed state and separate current-activation predicate.
The serialized Ready/Enter event tests and sets the durable entry flag
before dispatching signatures. Repetition, restart and late callbacks
therefore cannot repeat Enter or reset its entry time. Each core
message is authenticated and routed by the pair it signs; local tasks
use the same pair and anchor context. Waiting does not run the core.
Closing and Terminated records prevent a late event from reopening an
old program. These checks establish the stated isolation by induction
over the local event sequence.
\end{proof}

This is local authorization relative to a node's processed public
prefix. It does not claim that all nodes learn a finalized Force at the
same real time. A lagging node can briefly run an old program, but
\cref{prop:global-prefix} still excludes its closed-version output.
\Cref{prop:retirement-safety} checks the additional regional retirement
behavior. Neither proposition independently re-proves all of Part~I
or the baseline DATA/DRAIN and application-validation interfaces.

\subsection{Progress conditions and limits}
The regional engine keeps its original fair-schedule, exclusion and
uninterrupted-commit-window conditions. Normal handoff still needs
state transfer and authority, but no Espresso confirmation of its
terminal DATA when direct transfer succeeds. Espresso retrieval provides
the alternative path. Global output uses the Publish Layer and Espresso
service bounds, but also needs complete next results and successful
acceptance before the fixed deadlines. The transport bounds alone do not
establish that end-to-end rate.

A Force needs a recoverable retained anchor and an eligible successor.
Its startup budget starts at the public opening, not at a replica's
later Ready time. Fixed budgets therefore also assume timely startup
for some serving interval. An empty DATA can satisfy protocol delivery
without establishing user-transaction inclusion or censorship resistance.
The revised selector permits one eligible region to renew its own turn,
subject to its governance state and cooldown. This removes one selection
stall but does not establish that its next result beats a fixed cutoff.

Espresso's own failure prevents new global confirmation and automatic
crossings. A future-dated multisignature Force removes the need for a
second L1 round trip, but still waits for its announced Espresso
boundary and subsequent Ready. This document gives no unconditional
wall-clock bound on failover or withdrawal.

\subsection{What the selection repairs establish}
\label{sec:selection-repairs}
\begin{proposition}[Advancing normal selection and same-region renewal]
\label{prop:selection-repairs}
Along consecutive regular handoffs, the opening selection heights strictly
increase. If there are infinitely many such handoffs and eligibility is
a fixed nonempty set after some Espresso height, each member of that set
is selected repeatedly. Separately, a fixed pending boundary with no
eligible alternative can select its own source once that source is
eligible. Neither statement proves timely global acceptance.
\end{proposition}
\begin{proof}
For a regular successor,
\[
 q^{\rm open}_{\beta,r+1}
 =q_s\ge \operatorname{height}(\xi(T_{\beta,r}))
 >q^{\rm open}_{\beta,r}.
\]
These heights are integers, so infinitely many regular handoffs cannot
stay below any fixed historical height. Once past the eligibility-change
height, the cyclic search visits the fixed eligible set in order. If it
has one member, the fallback renews that member; otherwise the search
returns the next member before returning to the source. A finite cyclic
list therefore visits each eligible member repeatedly.

For the second statement, the historical search takes the least position
at which its selector returns a region. If no other region is eligible,
that selector returns the source exactly when the source itself is eligible.
An AUTO-imposed finite cooldown can expire under growing Espresso time;
this does not override OFF. A later effective Force may supersede the
pending boundary, which is why the statement fixes the boundary.
Selection creates a new pair; it neither runs Enter nor accepts DATA.
\end{proof}

The fixed-eligibility and infinite-handoff conditions specify this
proposition's scope; they are not added environmental assumptions for
the complete protocol. In particular, this proposition does not prove
that such a successful infinite sequence exists.

\subsection{Liveness targets not yet established}
\label{sec:liveness-counterexamples}
The complete target still includes continuing global progress with a
usable region, useful service for healthy eligible regions, and request
progress. These are stronger than selecting a region or keeping old
DATA recoverable. The independent v3.2 review gives the following
obstructions, which the interface and selector repairs do not remove.

\paragraph{A stronger slow-session trace.}
Use regions A and B, each with two correct unit-weight members and one
silent Byzantine member, $W=3,F=1,\Delta=1$. All authorities and state
are available. Later duties have quota one, startup target one,
$W_{\rm start}=6$ and $C_{\rm auto}=1/2$. No committee operation
interrupts them. Espresso timestamps equal real time, and complete
results meet $T_{\rm pub}=T_E=1$.

After an arbitrarily long finite successful prefix, open a duty at time
$a$. In v3.3 the next Espresso header can appear at $a+\varepsilon$,
for $0<\varepsilon<1/10$; both correct members pass Ready and Enter then.
The correct initializer sends NEWVIEW at $a+2+\varepsilon$, obtains
the other member's ANCHOR vote at $a+4+\varepsilon$, and proposes the
terminal DATA. It votes immediately; the other member votes at
$a+5+\varepsilon$. The two commit timers expire at
$a+6+\varepsilon$ and $a+7+\varepsilon$. The earliest complete FINAL
receipt is no earlier than the latter time when the Byzantine member
withholds its FINAL.

AUTO already closed this version at $a+6$, with accepted count zero.
Closure does not cancel those timers; the complete result later forms
and is published within its stated service bounds, but cannot enter
$G$. The successor opens at $a+6$. The prior source's cooldown has
expired before it is next needed at $a+12$. Repeat while alternating
each region's correct initializer identities. Every correct identity
actually enters, proposes DATA and commits infinitely often; this
trace does not exploit an unused nominal leader schedule. Strictly
advancing references and complete receipt multicast do not make a
complete result exist before its first honest commitments.

The finite-interval reading of \emph{some sustained service interval} is
satisfied by the successful prefix. Under that reading, $G$ then stops.
If the adjacent service prose instead requires sufficiently long
successful intervals for every future duty or an infinite tail, the
trace is excluded by that stronger interpretation. The frozen text
does not formalize that quantifier. Neither interpretation can be
silently substituted for the other in a liveness proof.

\paragraph{Global growth does not prove request progress.}
Give a correct admissible request $x$ to every correct node and put it
first in their queues. Every correct-led slow DATA above includes $x$.
Interleave fast sessions led by a Byzantine initializer which follows
the signing rules but proposes an empty DATA. With sufficiently small
actual message delays, after its required $2\Delta$ wait this session
forms its receipt at about $a+4$, and its EXPORT can be accepted before
$a+5<a+6$. Both regions can receive such successful sessions repeatedly.
The remaining turns rotate among all correct identities and follow the
slow trace. They really produce $x$, but only in closed suffixes.
No equivocation is needed, so eventual equivocator exclusion does not
remove the fast initializer. Empty DATA is valid and counts for delivery.

Thus $G$ and each region's accepted output can grow forever while $x$
never enters $G$. Pending-request preservation and finite older queues
hold. Correct leader fairness alone does not imply that some correct-led
output wins the global deadline. A stronger acceptance-service reading
must explicitly cover the required correct output or requests, rather
than merely some other successful EXPORT.

\paragraph{What the earlier examples do and do not show.}
The old constant-$q_0$ A/B cycle could keep a later-ON C out of historical
selection. Its status as a counterexample to all progress premises
depended on whether leader fairness already quantified over C's leaders.
The new strict-reference rule removes that protocol-level cycle without
claiming it was valid under every fairness interpretation.
Likewise, v3.2's sole-eligible-region stall violated the full premise
that an unassigned boundary eventually has another eligible region.
The fallback removes that rule-level exclusion; it does not prove the
stronger one-live-region goal against continual preemption.

The two service quantifiers need an exact interpretation, and fixed
cutoffs still need a protocol repair that does not repeatedly invalidate
all useful work. Increasing a constant alone does not bound unbounded
state preparation or recovery-authority waits. An architecture with
protected work or deferred effective Force would change cutoff semantics
and must say so explicitly. This version does not make that change or
claim the complete goal achieved.

\Needspace{18\baselineskip}
\subsection{Open design and implementation questions}
\begin{enumerate}[label=(\arabic*)]
\item Pin the concrete Espresso finality, timestamp and namespace-proof
profile, including clock lead and history retention.
\item Fix registry/control-proof encodings and the certifier's
complete-control-version acceptance rule.
\item Choose window and cooldown parameters from explicit service
assumptions, and address the service-window failures in the preceding section.
\item Specify native execution and signing continuation at every
allowed global anchor, including partial-EVM state if supported.
\item Independently verify the retirement extension and its conditional
regional safety argument; do not reuse the old universal-QC lemma after
the signing gate closes.
\item Complete staking retirement and show that the Publish Layer
implementation meets its service bound. A missed delivery window alone
is not a proof of proposer equivocation.
\end{enumerate}

\section{Implementing the Publish Layer}
\label{sec:publish-implementation}
The preceding protocols use a delivery-service assumption. The following
leader-first design explains how existing participants can implement it.
It introduces no new consensus round and no exclusive publisher role.
The implementation must meet $T_{\rm pub}$; the sketch alone is not a
proof of that bound.

\subsection{Leader first, with permissionless backup}
The proposing regional leader is the preferred sender. Once it has the
complete finalized result, it builds and submits an EXPORT immediately.
Other replicas continue collecting FINAL messages, retaining payloads and
making complete results available to the service. Commitment and
Completion do not wait for a publication acknowledgement.

Another replica or an external data verifier may submit the same result
without the leader's permission. An implementation may give the leader a
short first opportunity, but the backup delay, data retrieval and retries
must fit within $T_{\rm pub}$ from complete-result availability. A private
announcement that the leader has sent something is not enough to suppress
backup indefinitely. The service checks submission and Espresso evidence.

The fallback sender uses the original DATA and receipt. It cannot replace
the DATA, rewrite its session pair or create FINAL on behalf of the region.
The service guarantees delivery through an available path; no particular
sender is guaranteed to win the submission race.

\subsection{Fee sponsorship}
A backup sender need not hold tokens or pay from its own balance. It may
use a funded relay, a prepaid submission account, or another sponsored
route that can submit to Espresso. The route carries the verified result;
it does not acquire authority to change its content.

Someone still supplies the fee resources. To justify the publication
bound, a complete result covered by the service assumption must have access to a funded,
working route even when the original leader is unavailable. A sender can
try another route if one refuses or fails. The concrete sponsorship and
admission mechanism is an implementation choice; this protocol assumes
no Espresso application smart contract for it.

\subsection{Duplicates, retries and acceptance}
Different senders may carry the same result in different EXPORT records.
The DATA's producing pair $(\beta,r)$ and block hash identify the output;
receipt encoding, signer set and sender do not create a second output.
Only a valid next interval advances $G$ and its window count.

After submission, service workers track the Espresso record and replay
its acceptance. If another submission has already advanced $G$, they
rebuild an EXPORT for the remaining continuous suffix. DATA already in
$G$ is complete. A result rejected because of a gap must be resubmitted
after its predecessors; an old record is not accepted retroactively.
A closed version never regains permission through a retry, although its
record can still serve as historical regional evidence.

The two service bounds establish delivery and visibility of a submitted
record. They do not by themselves bound a sequence of stale-base retries.
Progress therefore still needs a serving interval in which complete next
results are accepted before the delivery deadlines, as stated in
\cref{sec:progress}. Regional production and state transfer must also fit
that interval.

\subsection{Retention and incentives}
Before finalized Espresso inclusion, the holders of a complete result
keep its payload and evidence available to backup senders. Afterward,
Espresso's retained, authenticated history is the common retrieval path.
Direct transfer can still accelerate regular handoff and state recovery.
Retirement and checkpoint rules must preserve every required copy.

Submission rewards and reimbursement may fund this service. They must not
count duplicate DATA as new progress. A missed delivery window alone does
not prove that the original leader, a verifier or a sponsor equivocated.
This document does not derive a slashing rule from that timeout.

\Needspace{14\baselineskip}
\begingroup\small
\begin{thebibliography}{9}
\bibitem{part1}
Sync HotStuff RR, Part I: Core Protocol. Design document, 13 October 2026.
Companion source: \href{run:sync_hotstuff_rr_part1_core.tex}{Part I core protocol}.
\bibitem{sync}
I. Abraham, D. Malkhi, K. Nayak, L. Ren, and M. Yin.
\emph{Sync HotStuff: Simple and Practical Synchronous State Machine Replication}.
IEEE Symposium on Security and Privacy, 2020, pp. 106--118.
\href{https://conferences.computer.org/sp/pdfs/sp/2020/349700a106.pdf}{Published paper}.
\bibitem{certification}
MegaETH Specification, PR~8, \emph{Block Certification}.
Commit \texttt{aedc2826663b}.
\href{https://github.com/megaeth-labs/specification/tree/aedc2826663b134ee25d26bea17e70c6d23d321c/specs/block-certification}{Pinned specification}.
Its Espresso and certifier verification interfaces are a reference;
the rotation and window rules in this document are new.
\end{thebibliography}
\endgroup
\end{document}
